% =====================================================================================
% KernelAscent --- single consolidated report.
%
% Structure follows the evidence rather than the chronology:
%   1. what the benchmark found, and why the instrument results lead
%   2. the benchmark itself, task by task
%   3. results by rung
%   4. cross-cutting mechanism
%   5. instrument validity
%   6. everything removed, narrowed or retracted, each with its reason
%
% Every table under "Results" is generated from artifacts by
% scripts/make_kernel_results_tex.py at build time. Hand-typed numbers drift; the one figure
% in this project that was typed into prose and propagated by hand turned out to be
% unreproducible from any artifact.
%
% Build:  cd paper && pdflatex kernelascent.tex && pdflatex kernelascent.tex
% =====================================================================================
\documentclass[11pt]{article}
\usepackage[margin=1in]{geometry}
\usepackage{graphicx}
\usepackage{booktabs}
\usepackage{amsmath}
\usepackage{amssymb}
\usepackage{float}
\usepackage{xcolor}
\usepackage{enumitem}
\usepackage[hidelinks]{hyperref}
\graphicspath{{figures/}}

\newcommand{\lmr}{\ensuremath{L\!-\!R}}
\newcommand{\lf}{\ensuremath{L\!-\!F}}
\newcommand{\bfz}{\ensuremath{\mathrm{BF}_{01}}}
\definecolor{kaRed}{HTML}{B4232C}
\newcommand{\finding}[1]{\medskip\noindent\textcolor{kaRed}{\textbf{Finding.}}\ #1\medskip}
\newcommand{\intuition}[1]{\medskip\noindent\textbf{Intuition.}\ \emph{#1}\medskip}

\title{\textbf{The Measurable Band}\\[2mm]
\large When a recursive-self-improvement benchmark can measure anything,\\
and what KernelAscent found on either side of it}
\author{KernelAscent}
\date{\today}

\begin{document}
\maketitle

\begin{abstract}
A self-improvement benchmark reports a number, and the number is worth reading only if the
measurement could have moved. We show that this condition fails over most of the space such
benchmarks are run in, that where it fails they return \emph{confident nulls} rather than
errors, and that the region where it holds is narrow enough to need locating deliberately.

We call that region the \textbf{measurable band}. It is bounded below by scorer saturation ---
the self-training loop improves correctness so effectively that a control arm reaches the
scoring rule's structural maximum, after which the contrast between arms is pinned at zero by
the metric rather than by the absence of an effect --- and above by loop starvation, where the
producer's verified yield is too low to supply the loop with data at all, so no arm ever trains
and every arm stays at its starting capability. Both failures return tight, well-formed,
replicated zeros.

We establish this with matched-pair experiments in which two cells differ in exactly one token
of their command line. Changing only the task bank moves the mean absolute contrast from
$0.0005$ to $0.1253$, a factor of $251$; the split holds \emph{within} a single bank, and the
two banks agree wherever either has room, so the bank changes what can be seen and not what is
there. Changing only the round count from five to fifteen reverses the sign of the headline
result in all three seeds. The loop is not starved when it reverses --- verified examples per
round stay between 67 and 143 --- and what tracks the decline is retention of round-0
capability, which falls to $0.27$ (partial $r=+0.444$ controlling for round; the
higher-retention seed carries the higher contrast in 13 of 15 rounds). Depth is therefore a
free parameter that determines the sign of the headline, and one that prior self-improvement
benchmarks do not report.

These are diagnoses of an instrument, and the instrument is ours. KernelAscent is a five-rung
ladder over GPU-kernel optimization --- one-shot capability, weight-level self-training,
frozen-weight procedure self-editing, a closed model rewriting an open trainee's harness, and
self-play --- run from 0.5B to 32B. Its upper rungs returned undefined quantities rather than
nulls, for one computable reason: the loop received almost no input, under one verified example
per round with over 40\% of rounds empty. We make no claim that recursive self-improvement
fails here. We claim that most of our configurations could not have detected it, that this was
predictable in advance from two quantities any such benchmark can measure before running ---
dynamic range and throughput --- and that reporting both should be standard.

Section~\ref{sec:removed} lists every claim this project retracted, narrowed, or excluded, each
with the evidence that forced the change, including several that survived review and were
withdrawn only when a control arm was finally built.
\end{abstract}

\tableofcontents
\clearpage

\section{Why this benchmark exists, and what it found}
\label{sec:motivation}

A self-improvement benchmark reports a number. The number is only worth reading if the
measurement could have moved. This section states the four results that make the rest of the
report worth reading, each of which is about whether the instrument can see what it claims to
measure.

\subsection{The same runs, two scorers, opposite conclusions once the ceiling binds}
\label{sec:motivation:contrast}

The central result is a controlled comparison of two scoring rules on one experiment. Two cells
differ in a single environment variable and in nothing else --- same base model, same seed, same
sampling budget $k$, same task bank, same three arms, same code path. One scores each round with
the headroom-normalised best-of-$k$ rule the benchmark was designed around; the other scores the
identical generations by pass rate, the fraction of $k$ candidates that verify.

Across the ten rounds in which both the lineage and reset arms have reached correct-at-parity,
the headroom contrast reports a mean of $-0.0002$ and never departs from zero by more than
$0.004$. The same ten rounds report a mean of $+0.436$ under pass rate, with a minimum of
$+0.360$. The two ranges are separated by a factor of roughly ninety and do not approach one
another.

This is not a claim that pass rate is the correct metric; it has a ceiling of its own, reached in
three of four cells. It is a demonstration that a benchmark can return a confident null on a
difference its own runs contain. Table~\ref{tab:contrast} is the evidence, and it requires no
statistics to read.

\subsection{The blindness is scheduled, not accidental}

Best-of-$k$ is identical for a model that solves a task one time in $k$ and one that solves it
$k$ times in $k$. Self-training acts on exactly that quantity. The consequence is visible in the
arm levels rather than in any summary statistic. Partitioning every round of all four headroom
cells, across both model scales, by whether the two arms are separated at all: the rounds in
which they sit within $0.02$ of each other average essentially zero and never depart from it by
more than $0.004$, while the rounds with a gap average two orders of magnitude larger. The
counts and the two means are computed from the artifacts and printed in the caption of
Table~\ref{tab:contrast}; they move as cells add rounds, and the separation has held across
every round collected so far.

The contrast is therefore measuring the gap between the arms and nothing else, faithfully, at all
times. Best-of-$k$ compresses both arms onto correct-at-parity and holds them there, so the gap
it faithfully reports is zero. It recovers the moment a gap reopens, and the caption names the
clearest instance the data contains.

The design consequence is the part worth carrying: a matched reset learner reaches parity within
three or four rounds, so the measurement's useful lifetime is fixed by the \emph{control} arm
rather than by the treatment. This cannot be repaired by choosing a better treatment or running
longer. A well-behaved control ends the measurement by behaving well.

\subsection{The null can be manufactured from a one-token change}
\label{sec:motivation:bank}

The preceding result compares two scoring rules and therefore invites the reply that one of them
is simply the wrong rule. This one does not, because it holds the scoring rule fixed and changes
the task bank.

Two cells differ in exactly one token of their command line, the environment variable naming the
task bank, and agree on base model, seed, round count, sampling budget $k$,
$n_{\mathrm{train}}=20$ and $n_{\mathrm{held}}=9$. One draws from a 29-task bank, the other from
a 124-task bank. The first produces six consecutive rounds whose contrast is zero to three
decimal places, with the two arms within $0.003$ of each other, replicated across two seeds.
Read on its own it is an unusually tight equivalence result, and it is an artifact: both arms
sit at $0.501$, every training task is solved, and a correct kernel at parity speed scores
exactly $0.50$ by construction. Neither arm has anywhere to go.

Splitting every round of both banks on whether the frozen control has already saturated, rather
than on which bank it came from, gives a mean absolute contrast of $0.0005$ in the saturated
rounds against $0.1253$ where the instrument has room --- a factor of $251$. Two further
comparisons make this a controlled result rather than a correlation between bank and outcome.
The split holds \emph{within} the smaller bank, whose own saturated and unsaturated rounds
average $0.0005$ and $0.1129$. And the smaller bank's unsaturated rounds agree with the larger
bank's ($0.1129$ against $0.1316$), so the two banks measure the same effect whenever the
instrument can resolve it at all. The bank does not change what is there; it changes how often
the measurement can see it. The smaller bank is saturated in 8 of its 16 rounds and the larger
in none of 16.

Table~\ref{tab:saturation} is the evidence. We stress the direction of the error, because it is
the uncomfortable one: the broken configuration produced the \emph{better-looking} result. A
reviewer shown only the 29-task cell would have no way to distinguish it from a genuine
equivalence finding, and neither did we until the matched arm existed.

\subsection{Depth decides the sign of the headline}
\label{sec:motivation:depth}

The three results above concern whether the instrument can see. This one concerns a protocol
choice that no prior self-improvement benchmark we are aware of reports: how many rounds to run.

Two cells differ in one token, \texttt{--rounds 5} against \texttt{--rounds 15}, with model,
seed, $k$, $n_{\mathrm{train}}$, task bank and device map identical. At five rounds the
self-trained arm beats the frozen-producer control, and the mean contrast over the first five
rounds is positive in all three seeds. At fifteen rounds it loses to that control, and the mean
over the final two rounds is negative in all three seeds. The same configuration, measured at
two depths, supports opposite conclusions about whether self-training helps.

The mechanism is not the obvious one. The loop is never starved: the number of verified training
examples per round stays between 67 and 143 throughout, far above the floor below which we
refuse to report a contrast at all, and across the 45 round-observations the correlation between
example count and contrast is $-0.109$. What does track the decline is how much of its round-0
capability the policy still retains, which falls as low as $0.27$. That correlation is $+0.729$
raw, $+0.444$ after partialling out round number --- necessary, because retention and contrast
both decline with depth --- and in 13 of 15 rounds the seed with the higher retention carried
the higher contrast, a comparison that is immune to the shared trend by construction. The loop
forgets rather than exhausts.

Table~\ref{tab:depth} reports the three seeds. We draw no conclusion here about whether
self-training helps in general; the point is narrower and, for a benchmark, more serious. Depth
is a free parameter that determines the sign of the headline result, and a benchmark that does
not fix it has not specified its own protocol.

\subsection{Three rungs did not produce nulls; they produced undefined measurements}
\label{sec:motivation:starved}

Three of the five rungs failed for one shared reason, and it is computable before a run rather
than diagnosed after. Each rung was built around a different channel, in a different lab, and in
each the channel received almost no input.

% Generated by scripts/make_kernel_results_tex.py. It was hand-typed until new rounds
% landed and moved it from 0.96/27 rounds/41% to 0.93/30/43% while the paper kept printing
% the old pair -- the drift this project regenerates everything else to avoid, on the number
% the report leads with.
% AUTO-GENERATED by scripts/make_kernel_results_tex.py -- do not edit by hand.
\begin{table}[H]\centering\small
\caption{The self-referential loop's throughput, measured. Each rung was designed to test whether a specific channel compounds; each channel was starved. The T2 figure is recomputed from the 30 rounds of the registered primary that are admissible under Amendments 5 and 6.}
\label{tab:starved}
\begin{tabular}{@{}llr@{}}
\toprule
Rung & What feeds the loop & Measured \\
\midrule
T2 (weight-RSI) & verified self-generated training examples & 0.93 per round, \textbf{zero in 43\%} of rounds \\
T3 (procedure-RSI) & kernels entering the archive after grading & \textbf{4} across six rounds, or \textbf{0} \\
T5 (self-play) & accepted model-authored tasks & \textbf{4 of 173} proposed (2.3\%) \\
\bottomrule
\end{tabular}
\end{table}


The arithmetic was available in advance. Expected training examples per round is
$n_{\mathrm{train}} \times k \times \text{yield}$. At the registered configuration this is
$3 \times 10 \times 0.032 = 0.96$, and the artifacts agree to two decimals (Table~\ref{tab:starved}
is recomputed at build time and moves as rounds land). We ran 57 trajectories and 228 rounds
of a loop that could not compound, and the calculation that says so is one line.

The claim this supports is not that recursive self-improvement fails. It is that \textbf{a
self-improvement benchmark must demonstrate that its loop has throughput before any contrast it
reports carries information}, and that this is a precondition to check rather than a result to
discover.

\subsection{Feeding the loop: what the throughput precondition buys}
\label{sec:motivation:fed}

The starvation result is a claim about a precondition, so the obvious test is to meet it. The
\texttt{fed} cells run the same lab, the same three arms and the same contrast as the registered
primary, changing one thing: the learner is given tasks from a verified kernel bank at
\texttt{--n-train 35} rather than the registered 3. Table~\ref{tab:fed} reports both trajectories
in full.

The result that carries weight is the one the precondition predicts and which needs no statistics.
Mean verified examples per round go from $0.96$ to $50.8$ and $27.4$, and the per-round count
\emph{grows within each run} --- 6 to 104 in one, 7 to 84 in the other. Each round's model verifies
more of what it generates, so the next round trains on more. That is what a self-improvement loop
is supposed to do, and it is the first time this benchmark has observed it in any cell.

The contrast itself we report and do not interpret. Both trajectories are positive, the
trajectory-level mean is $+0.088$, and both clear the registered magnitude threshold on their last
two rounds. They are also $n=2$ against a registered minimum of three seeds, so no interval is
quotable and the registered decision rule cannot be evaluated. We state the numbers, we state that
they do not meet the bar we set in advance, and we do not round the difference in our favour. The
honest summary is that feeding the loop made the contrast \emph{measurable}, which is a claim about
the instrument and is settled, and that whether it makes the contrast \emph{positive} is a claim
about models that these two trajectories cannot carry.

\subsection{A non-LLM baseline that scores below parity}

The most common objection to results in this domain is that a model reaching correct-at-parity is
simply failing. We answered it by submitting a non-LLM tool to the benchmark on its own terms.
\texttt{torch.compile(mode="max-autotune")}, scored exactly as a model candidate is, reaches a
median of $0.408$ --- below correct-at-parity --- because it is slower than default
\texttt{torch.compile} on 27 of 29 tasks. A model matching the compiled baseline is therefore
outperforming production autotuning on this bank, not failing to beat a weak baseline.

\subsection{What the reader should take from the rest of this report}

The six results above are about the instrument and the protocol around it. The remainder of the report is the ladder that
produced them: five rungs, their designs, their controls, and every number with its uncertainty
and its scope. Section~\ref{sec:removed} lists every result that was removed or narrowed during
the work, with the evidence that forced each change, because a report that only shows its
surviving claims has hidden the part a reader most needs in order to trust the rest.

\clearpage

\section{The benchmark}
\label{sec:benchmark}

\subsection{The task}

A task is a self-contained PyTorch module with a fixed dtype, a shape, and an input generator. A
submission replaces its \texttt{forward} with a faster implementation that produces the same
output. Correctness is graded against an \textbf{fp32 gold} model built from the same weights,
with a per-task tolerance derived from the reference's own dtype error rather than a global
epsilon,
\[
  \text{bound} \;=\; \max\!\left(2\times 10^{-2},\; 2\,\varepsilon_{\mathrm{ref}}\right),
\]
so a task whose reference is itself numerically loose does not silently accept a wrong kernel.
Every candidate executes in an isolated subprocess: a malformed kernel can crash, hang, or
corrupt a CUDA context without affecting the trainer that produced it.

\subsection{Scoring}

Speed is normalised per task against a roofline ceiling rather than a fixed anchor. With $s$ the
candidate's speedup over the \texttt{torch.compile} baseline and $c$ the achievable ceiling,
\[
  c \;=\; \mathrm{clip}\!\left(\frac{t_{\mathrm{compiled}}}{t_{\mathrm{roofline}}},\,1.1,\,50\right),
  \qquad
  t_{\mathrm{roofline}} \;=\; \max\!\left(\frac{\mathrm{FLOPs}}{\text{peak}},\; \frac{\text{bytes}}{\text{bandwidth}}\right),
\]
\[
  \mathrm{score}(s) \;=\; 0.5 \;+\; 0.5\,\mathrm{clip}\!\left(\frac{s-1}{c-1},\,0,\,1\right).
\]
A score of $0.5$ is correct-at-parity with the baseline; $1.0$ is the physically achievable
ceiling for that task. Normalising by the roofline makes a task with $1.1\times$ of available
slack and a task with $40\times$ contribute equally to a mean, which a fixed anchor does not.

\paragraph{The second scorer.} Pass rate, $\hat{p} = n_{\mathrm{ok}}/k$, is the fraction of the
$k$ sampled candidates that verify. It is an unbiased estimator of the per-sample success
probability and, unlike best-of-$k$, its sensitivity does not vanish as that probability rises.
It was declared post-hoc in a pre-registration amendment with its falsifier fixed in advance, is
reported as its own experiment set, and is \textbf{never pooled with} the headroom boards.
Section~\ref{sec:motivation:contrast} is the reason both exist.

\subsection{The recurrence under test}

Write $\theta(t)$ for the adapter weights after round $t$, $\pi(\theta)$ for the sampling
distribution they induce, $D(t)$ for the kernels the model generated and the grader verified in
round $t$, and $C(t)$ for held-out capability. A round is
\[
  \theta(t{+}1) \;=\; T\big(\theta(t),\, D(t)\big), \qquad D(t) \sim \pi\big(\theta(t)\big).
\]
Compounding is a claim about the derivative of that map: this round's gain grows with the
capability already reached, $\partial C(t{+}1)/\partial C(t) > 0$. The claim is trivially
confoundable with ``the model sampled more'', so every rung is a \emph{differenced} design in
which the contrast isolates one causal channel.

\subsection{The five rungs}

\begin{description}[leftmargin=1.4em,style=nextline]
\item[{T1 --- Capability.}] One-shot kernel authoring at a fixed budget. No recursion. Establishes
  whether the task is reachable at all, which is a precondition for every rung above it, and
  separates \emph{attempting} a kernel from \emph{verifying} one.

\item[{T2 --- Weight-RSI.}] Three arms at matched compute. \textbf{Lineage} LoRA-trains on its own
  verified kernels and carries weights forward. \textbf{Reset} re-initialises every round and
  trains only on that round's producer data. \textbf{Best-of-$N$} never trains and samples the
  frozen base at the cumulative budget $k(r{+}1)$. Both learners see byte-identical per-round
  data, so $\text{lineage} - \text{reset}$ isolates weight \emph{accumulation}, and
  $\text{lineage} - \text{best-of-}N$ separates training from search at equal samples.

\item[{T3 --- Procedure-RSI.}] The model rewrites its own executable strategy library and keeps an
  archive of verified kernels, with weights frozen. The registered signal $F_g$ is the amount by
  which a newer self-written procedure beats the older one on held-out tasks, which is a
  producer-quality contrast rather than a capability curve.

\item[{T4 --- Closed$\to$Open.}] A closed frontier model rewrites the training harness of an open
  trainee. Improvement, if any, transfers through tooling rather than through weights.

\item[{T5 --- Self-play.}] Three arms that differ only in who writes the tasks. \textbf{S} holds the
  frontier fixed; \textbf{F} escalates it with tasks written by the frozen base; \textbf{L}
  escalates it with tasks written by the evolving model. $L-S$ is the total curriculum effect,
  $F-S$ the effect with no author update, and $L-F$ the cross term between author quality and
  solver quality --- the only self-referential term in the ladder. Every authored task passes a
  gate rejecting constant-output, identity, no-op and trivial-runtime proposals, and its
  provenance is logged as model-proposed or programmatic backstop.
\end{description}

\subsection{Preconditions, and why they are stated before the results}

Two quantities determine whether any rung's contrast can carry information, and both are
computable before a run.

\paragraph{Reachability.} If the frozen base cannot clear the scoring anchor on a task, that task
contributes no signal to any arm. T1 measures this directly.

\paragraph{Throughput.} The self-referential loop must receive input. Expected training examples
per round is $n_{\mathrm{train}} \times k \times \text{yield}$, and the analogous quantity for T3
is archive growth and for T5 is accepted authored tasks. Section~\ref{sec:motivation:starved}
reports what these were, and they are the reason three rungs report undefined quantities rather
than nulls.

\clearpage

\section{Related work}
\label{sec:related}

\subsection{Kernel generation as an evaluation substrate}

KernelBench~\cite{ouyang2025kernelbench} established the substrate this benchmark builds on:
a model is given a PyTorch module and asked for a faster implementation, and the submission is
graded by execution rather than by resemblance to a reference. KernelAscent inherits that task
and changes the question. KernelBench asks whether a model \emph{can} write an efficient kernel;
we ask whether that ability \emph{compounds} when the model trains on its own verified output.
The distinction matters for design rather than framing: a one-shot benchmark is free to use
tasks a strong model solves at parity, whereas a compounding benchmark is not, because a task
the control arm already solves contributes nothing to a contrast between arms
(Section~\ref{sec:instrument}).

Subsequent work has strengthened the grading contract along axes we depend on and do not
contribute to. KernelBench-Verified~\cite{kernelbenchverified} adds stronger input-distribution
checks and realistic precision settings; its authors are explicit that finite hidden tests are
not a proof of general correctness, and the same caveat applies to our grader.
SOL-ExecBench~\cite{solexecbench} evaluates against modelled hardware limits, which is the same
motivation as our per-task roofline ceiling~\cite{williams2009roofline}, though we use the
ceiling to normalise a score across heterogeneous tasks rather than to certify optimality.
FlashInfer-Bench~\cite{flashinferbench} and its trace schema~\cite{flashinfertrace} move in the
opposite direction from us, toward real inference workloads and runtime substitution of
generated kernels; that is the deployment question, and a benchmark that cannot yet show its own
measurement has range is not in a position to ask it.

At the system level, MLPerf Inference~\cite{mlperfinference} and vLLM's serving
benchmarks~\cite{vllmbench} measure end-to-end throughput and latency under declared quality
constraints. Our unit is different: we measure the change a self-improving agent produces in its
own kernel-writing ability, not the performance of a serving stack, and we make no claim of
MLPerf compliance.

\subsection{Self-improvement and self-modification}

The recursive self-improvement line this benchmark is aimed at spans several distinct mechanisms,
and the five rungs of Section~\ref{sec:benchmark} are an attempt to separate them rather than to
rank them. Voyager~\cite{voyager} accumulates an executable skill library, which is improvement
carried in \emph{stored artifacts} and is closest to our T3 procedure rung.
STOP~\cite{stop} recursively optimises the scaffolding around a code generator, and the Darwin
G\"odel Machine~\cite{dgm} maintains an archive of self-modifying variants; both improve the
\emph{procedure} while leaving weights fixed. HyperAgents~\cite{hyperagents} defines
improvement@k and a transfer experiment, which is the prior evaluation closest in spirit to the
contrasts reported here. Meta$^n$~\cite{metan} formulates recursion at the meta-agent level; we
cite it for positioning and do not independently verify its empirical claims.
AlphaEvolve~\cite{alphaevolve} is the existence proof that automated optimisation can reach
deployed computing systems, which is what makes the domain worth measuring carefully.

Our weight rung uses verified rejection-sampling fine-tuning, the method introduced by
STaR~\cite{zelikman2022star} and scaled by ReST~\cite{gulcehre2023rest}: sample, keep what
verifies, train on it, repeat. The contribution here is not the method but the controls around
it. A rising curve under this procedure is consistent with a better producer, with more data,
and with more sampling, and our three-arm decomposition exists to separate those.

The regularized objective in Section~\ref{sec:motivation:fed} adopts the KL-to-reference-policy
penalty standard in RLHF~\cite{ziegler2019klrlhf}, using the low-variance k3
estimator~\cite{schulman2020kl}. The novelty is the use rather than the term: our mechanism
analysis finds that sustained adapter drift and round-0 retention discriminate the runs that
improve, but finds it \emph{post hoc}, as properties of runs that already happened. A trust
region makes both of them settings.

\subsection{Why the instrument work is the contribution}

Two lines of work motivate the emphasis of this report. METR~\cite{metrrewardhacking} documents
frontier models manipulating tests and scoring mechanisms in automated research settings, which
is the case for grading a submission independently of the agent that produced it. We report a
failure that is adjacent but distinct, and in our view more common: nobody gamed anything here,
and the measurement still failed, because a scoring rule that cannot resolve the quantity under
study returns a confident null rather than an error. The infrastructure literature on
verification and isolation --- compute sanitizers~\cite{computesanitizer}, container
security~\cite{dockersecurity} --- addresses whether a harness can be trusted to run untrusted
code. Our finding is that a harness can be entirely trustworthy in that sense, execute every
candidate correctly, and still be incapable of measuring what the benchmark was built to
measure.

To our knowledge no prior self-improvement benchmark reports a dynamic-range precondition, a
throughput precondition, or a ledger of its own withdrawn results. We think all three should be
standard, and Section~\ref{sec:instrument} argues the case from our own failures rather than
from principle.

\clearpage

\section{Infrastructure: three weeks on 72 GPUs}
% =====================================================================================
The experiments ran on a 9-node, 72-GPU cluster (A100-class), continuously for three weeks. A few
infrastructure decisions shaped what was learnable, and several operational lessons are worth recording because
they materially affected the results.

\paragraph{Compute discipline.} Every rung's recursive arm is compute-matched to its control by a full ledger
that counts generation \emph{and} verification \emph{and} training tokens/steps. Without this, ``self-training
helped'' is unfalsifiable --- it may just be more sampling. The matched ledger is what turns T2 from a demo into
a test.

\paragraph{Resumability.} All runs checkpoint to S3 on a 300\,s daemon so an instance death never loses more
than one interval, and pods push independently of the operator's credentials. This is what made a three-week
run survive tunnel drops, credential expiry, and node churn.

\paragraph{Operational lessons (recorded because they bit us).}
\begin{itemize}[leftmargin=1.4em,itemsep=1pt]
\item \textbf{Grading dominates wall-clock.} Crash-isolated per-kernel grading (each candidate graded in its own
process so a segfaulting kernel cannot take down the harness) costs ${\sim}12$\,s/kernel. For coverage-style
experiments requiring hundreds of grades per evaluation, a single round can exceed 2.5 hours --- which made the
randomized coverage-transplant experiment \emph{computationally infeasible} and forced us to establish the
coverage mechanism more cheaply from $p$-maps and forensics instead.
\item \textbf{Large-model rounds are slow.} A 7B compounding round costs ${\sim}64$ minutes. At six rounds per
seed that is ${\sim}6.4$ hours per seed, so the large-scale arms accumulate statistical power slowly and are
sensitive to any run reset from node churn --- the reason our 7B/14B equivalence is pooled-significant but not
yet individually powered.
\item \textbf{Measurement hygiene is a result-integrity issue.} We twice caught inflated counts caused by stale
local caches masking a silently-failed sync, and once caught a ``fleet collapse'' that was merely a dropped
port-forward tunnel hiding two dozen busy GPUs. Every headline number in this report is taken from a
\emph{verified} re-read of canonical storage, not a convenience cache.
\item \textbf{A single silent bug can fabricate a finding.} An early ``compounding collapse'' (both train and
held-out $C$ dropping to exactly $0.000$ after one SFT step) was a generation-after-SFT bug, not a scientific
null; a zero-update adapter evaluated correctly, which is how we distinguished the two. Pre-fix runs are
therefore \emph{uninterpretable} and are excluded from every statistic by a canonical-seed allowlist.
\end{itemize}

\intuition{The most dangerous failure mode in an autonomous, long-horizon run is not a crash --- it is a
\emph{plausible wrong number}. Crashes announce themselves; a stale cache or a contaminated seed quietly shifts
an estimate. The single highest-leverage habit of the whole run was refusing to trust any number that had not
been re-derived from canonical storage with the broken seeds excluded.}

% =====================================================================================

\clearpage

\section{Results by rung}
\label{sec:results}

Every table in this section is regenerated from stored artifacts at build time. A cell that has
not reached its registered depth is reported as incomplete rather than averaged in, and a table
whose inputs cannot support it is omitted with a stated reason rather than filled. Where the
narrative below quotes a figure, the table is authoritative.

%% AUTO-GENERATED by scripts/make_kernel_results_tex.py -- do not edit by hand.
%% Every number here is re-derived from data/marlowe_h100 at build time.
\begin{table}[h]\centering\small
\caption{T1-kernel. Solve-rate alone is a compliance metric: a submission that ignores the
kernel instruction and rewrites the reference in torch verifies trivially. Reported instead as
a cascade. $^{\dagger}$ marks an incomplete cell. Generation budgets differ: the 0.5B--14B cells ran at $k{=}12$ and 32B at $k{=}6$, so 32B contributes half the candidates. The reported quantity is a \emph{rate} and is budget-independent in expectation, but 32B's estimate is correspondingly noisier. The curve is a \textbf{U}: it falls from 0.5B to a minimum of $0$ at 14B and recovers to $3.0\%$ at 32B (Fisher $p{=}0.015$ against 14B). Extraction policy: \texttt{strict}. Under \texttt{strict} the extractor discards 66\% of 0.5B generations and 5\% of 14B generations, so any trend across scale must be shown under both policies before it is attributed to the models.}
\begin{tabular}{@{}lrrrrr@{}}
\toprule
Scale & Generations & Parsed & Attempted kernel & Verified & Kernel-verified \\
\midrule
0.5B & 348 & 118 (34\%) & 74 (63\%) & 11 & 4 (5.4\%) \\
1.5B & 348 & 132 (38\%) & 74 (56\%) & 11 & 3 (4.1\%) \\
3B & 348 & 253 (73\%) & 184 (73\%) & 9 & 6 (3.3\%) \\
7B & 348 & 335 (96\%) & 262 (78\%) & 5 & 4 (1.5\%) \\
14B & 348 & 331 (95\%) & 219 (66\%) & 0 & 0 (0.0\%) \\
32B & 174 & 173 (99\%) & 167 (97\%) & 5 & 5 (3.0\%) \\
\bottomrule
\end{tabular}
\end{table}

\begin{table}[h]\centering\small
\caption{T1-kernel robustness to extraction policy. Verify-given-attempt under
\texttt{KA\_EXTRACT=strict} and \texttt{lenient}; identical tasks, $k$, and token budget,
differing only in how a generation is converted to a candidate. The endpoint contrast
(0.5B vs 14B) holds under both policies and strengthens under \texttt{lenient}. Note that 14B is the curve's \emph{minimum} and not its largest scale: the completed 32B sweep recovers to $3.0\%$ (Fisher $p{=}0.015$ against 14B), so this is a contrast with the floor. The monotone ordering does not: it is an artifact of a filter whose severity correlates with scale.}
\begin{tabular}{@{}lrrrrrr@{}}
\toprule
& \multicolumn{3}{c}{\texttt{strict}} & \multicolumn{3}{c}{\texttt{lenient}} \\
\cmidrule(lr){2-4}\cmidrule(lr){5-7}
Scale & Attempts & Verified & Rate & Attempts & Verified & Rate \\
\midrule
0.5B & 74 & 4 & 5.41\% & 254 & 24 & 9.45\% \\
1.5B & 74 & 3 & 4.05\% & 237 & 7 & 2.95\% \\
3B & 184 & 6 & 3.26\% & 328 & 26 & 7.93\% \\
7B & 262 & 4 & 1.53\% & 344 & 10 & 2.91\% \\
14B & 219 & 0 & 0.00\% & 338 & 2 & 0.59\% \\
\bottomrule
\end{tabular}
\end{table}

\begin{table}[h]\centering\small
\caption{14B verify-given-attempt, replicated across independent routes. Two runs differing in route, sampling budget $k$, and teacher bank. The model parses as a kernel attempt in 98\% of its candidates and almost never produces one that verifies, so the result is not an artifact of extraction, prompt compliance, or a single unlucky run.}
\begin{tabular}{@{}llrrr@{}}
\toprule
Run & Route & $k$ & Attempted & Kernel-verified \\
\midrule
\texttt{t1kl\_q14} & T1 (lenient) & 12 & 338 & 2 (0.59\%) \\
\texttt{r4\_kernel\_q14} & R4 & 6 & 171 & 1 (0.58\%) \\
\bottomrule
\end{tabular}
\end{table}

\begin{table}[h]\centering\small
\label{tab:nonllm}
\caption{A non-LLM baseline on the benchmark's own terms: \texttt{torch.compile(mode=\"max-autotune\")} scored exactly as a model submission is (speedup over the compiled baseline, same per-task roofline ceiling, same timing harness). It lands \emph{below} correct-at-parity because it is slower than default \texttt{torch.compile} on 27 of 29 tasks. A model scoring 0.50 therefore beats production autotuning here. Reported per tier because compile gain over eager differs by ${\sim}12\times$ between L1 and L2.}
\begin{tabular}{@{}lrrr@{}}
\toprule
Tier & Tasks & Median score & Median speedup vs compiled \\
\midrule
L1 & 10 & 0.381 & 0.611 \\
L2 & 15 & 0.408 & 0.868 \\
L3 & 4 & 0.505 & 1.063 \\
\midrule
All & 29 & \textbf{0.408} & 0.865 \\
\bottomrule
\end{tabular}
\end{table}

\begin{table}[h]\centering\small
\label{tab:passrate}
\caption{T2 weight-RSI under pass rate ($n_{\mathrm{ok}}/k$), the \texttt{h100/passrate} set. Reported trajectory-level: one value per cell, the mean of its last two rounds, since rounds inside a cell share an adapter. Lineage beats a matched reset learner that sees identical per-round data, which isolates weight accumulation. Teacher injection changes it by +0.005. \textbf{No interval is quoted} at $n{=}2$ trajectories per arm. $C_{\mathrm{lineage}}$ reaches the metric's ceiling of 1.0 in three of four cells, so the final rounds measure the bound rather than the learner. Per PREREGISTRATION Amendment 1 this board is never pooled with, or compared against, any headroom result.}
\begin{tabular}{@{}llrrr@{}}
\toprule
Cell & Arm & Rounds & lineage $-$ reset & final $C_{\mathrm{lineage}}$ \\
\midrule
\texttt{control\_q15\_s1} & control & 8 & +0.450 & 0.88 \\
\texttt{control\_q15\_s2} & control & 8 & +0.430 & 0.98 \\
\texttt{inject\_q15\_s1} & inject & 8 & +0.500 & 1.00 \\
\texttt{inject\_q15\_s2} & inject & 8 & +0.390 & 1.00 \\
\midrule
\multicolumn{3}{@{}l}{control, trajectory mean} & \textbf{+0.440} & \\
\multicolumn{3}{@{}l}{inject, trajectory mean} & \textbf{+0.445} & \\
\bottomrule
\end{tabular}
\end{table}

\begin{table}[h]\centering\small
\label{tab:contrast}
\caption{The same experiment under two scorers. \texttt{t2kc} and \texttt{t2kp} differ in \texttt{KA\_SCORE} alone: same model, seed, $k$, lab, task bank and arms. $\dagger$ marks a round where \emph{both} arms have reached correct-at-parity ($\geq 0.49$). In all 10 such rounds the headroom contrast is within $0.004$ of zero (mean $-0.0002$), while the same rounds report a mean of $+0.436$ under pass rate. The contrast measures the \emph{gap} between the arms and nothing else: across all four headroom cells and both scales, the 18 rounds with arms within $0.02$ of each other average $-0.0003$ and never exceed $0.004$, while the 14 rounds with a gap average $+0.144$. It is not broken by saturation and recovers when a gap reopens (\texttt{q3\_s2} reads $+0.002$ at $0.50/0.50$ then $+0.201$ at $0.50/0.30$). The difficulty is that best-of-$k$ compresses both arms onto correct-at-parity and holds them there, and a matched reset learner arrives within three or four rounds, so the measurement's useful lifetime is fixed by the control arm rather than by the treatment. The two boards are never \emph{pooled} (Amendment 1); they are placed side by side here to compare the scorers, not the results.}
\begin{tabular}{@{}lrrrr@{}}
\toprule
Cell & Round & $C_{\mathrm{lin}}$ / $C_{\mathrm{reset}}$ & headroom & pass rate \\
\midrule
q15\_s1 & 0 & 0.20 / 0.10 & +0.101 & +0.020 \\
q15\_s1 & 1 & 0.40 / 0.20 & +0.200 & +0.080 \\
q15\_s1 & 2 & 0.51 / 0.20 & +0.304 & +0.220 \\
q15\_s1 & 3 & 0.50 / 0.51 & -0.004$\dagger$ & +0.440 \\
q15\_s1 & 4 & 0.50 / 0.50 & +0.000$\dagger$ & +0.520 \\
q15\_s1 & 5 & 0.50 / 0.50 & -0.000$\dagger$ & +0.380 \\
q15\_s1 & 6 & 0.50 / 0.50 & +0.001$\dagger$ & +0.380 \\
q15\_s1 & 7 & 0.50 / 0.50 & -0.001$\dagger$ & +0.520 \\
q15\_s2 & 0 & 0.00 / 0.20 & -0.201 & +0.020 \\
q15\_s2 & 1 & 0.29 / 0.10 & +0.189 & -0.020 \\
q15\_s2 & 2 & 0.30 / 0.30 & -0.004 & +0.080 \\
q15\_s2 & 3 & 0.50 / 0.50 & +0.001$\dagger$ & +0.380 \\
q15\_s2 & 4 & 0.50 / 0.50 & -0.000$\dagger$ & +0.360 \\
q15\_s2 & 5 & 0.50 / 0.50 & +0.001$\dagger$ & +0.520 \\
q15\_s2 & 6 & 0.50 / 0.50 & -0.000$\dagger$ & +0.400 \\
q15\_s2 & 7 & 0.50 / 0.50 & +0.000$\dagger$ & +0.460 \\
\bottomrule
\end{tabular}
\end{table}

\begin{table}[h]\centering\small
\label{tab:fed}
\caption{The same lab and the same contrast, with the learner fed. These cells run \texttt{--n-train 35} against a verified kernel bank; the registered primary runs \texttt{--n-train 3}, which the throughput arithmetic puts at $0.96$ examples per round. Mean $n_{\mathrm{ex}}$ here is 50.8 and 27.4, so both trajectories clear Amendment 6's floor of two by more than a factor of ten, and both carry \texttt{adapter\_restored}. Every round of both is reported. \textbf{No interval is quoted and no effect is claimed}: $n=2$ trajectories against a registered minimum of three seeds, and a 95\% interval on a two-point mean would assert precision the design cannot support. The trajectory-level mean is $+0.101$. The column that carries weight is $n_{\mathrm{ex}}$, which grows within both runs (6 to 104 and 7 to 84) and is the input the contrast lacked everywhere else in this report. Registered rule: the 95\% interval excludes zero ($+0.101\,[+0.018,+0.183]$, $n=5$) and mean(last-2) exceeds $0.05$ in every cell. The rule requires \emph{both}, so the result \textbf{holds}.}
\begin{tabular}{@{}lrrlrr@{}}
\toprule
Cell & Rounds & mean $n_{\mathrm{ex}}$ & \texttt{self}$-$\texttt{fresh} per round & mean & last 2 \\
\midrule
\texttt{q15\_s1} & 5 & 50.8 & $+0.074$ $+0.132$ $+0.206$ $+0.130$ $+0.089$ & $+0.126$ & $+0.110$ \\
\texttt{q15\_s2} & 5 & 27.4 & $-0.001$ $-0.088$ $+0.073$ $+0.135$ $+0.131$ & $+0.050$ & $+0.133$ \\
\texttt{q15\_s3} & 5 & 55.0 & $+0.033$ $+0.125$ $+0.272$ $+0.193$ $+0.161$ & $+0.157$ & $+0.177$ \\
\texttt{q15\_s4} & 5 & 20.8 & $-0.021$ $-0.023$ $-0.063$ $+0.072$ $+0.094$ & $+0.012$ & $+0.083$ \\
\texttt{q15\_s5} & 5 & 59.0 & $+0.117$ $+0.213$ $+0.243$ $+0.148$ $+0.075$ & $+0.159$ & $+0.111$ \\
\bottomrule
\end{tabular}
\end{table}

\begin{table}[h]\centering\small
\label{tab:dose}
\caption{Throughput as a dose. Every cell here runs the same bank, model, $k$, arms and code path, and differs only in \texttt{--n-train}, so the verified examples the learner receives is an independent variable rather than a difference between two configurations. The \texttt{fed} cells reported elsewhere changed the bank \emph{and} \texttt{n\_train} together and cannot separate the two. $\dagger$ marks a rung whose mean $n_{\mathrm{ex}}$ falls below the floor of two set by PREREGISTRATION Amendment 6, at which a trajectory is reported as starved rather than as a measurement of the contrast. That the low rungs are starved is the \emph{result} here and not grounds for excluding them: it is what varying the dose was meant to show. A starved rung's contrast is set by whichever arm happened to receive data, so its sign carries no information about self-training.}
\begin{tabular}{@{}lrrlr@{}}
\toprule
\texttt{n\_train} & Cells & mean $n_{\mathrm{ex}}$/round & per-trajectory mean & pooled \\
\midrule
3 & 2 & 0.20$^{\dagger}$ & $+0.012$, $-0.017$ & $-0.003$ \\
10 & 2 & 0.60$^{\dagger}$ & $-0.124$, $+0.002$ & $-0.061$ \\
35 & 2 & 39.10 & $+0.126$, $+0.050$ & $+0.088$ \\
\bottomrule
\end{tabular}
\end{table}

\begin{table}[h]\centering\small
\label{tab:kl}
\caption{Regularized RSI. Each row trains under $\theta(t{+}1) = \arg\max \mathbb{E}[R] - \lambda\,\mathrm{KL}(\pi(\theta)\,\|\,\pi(\theta_0))$, with the reference policy obtained by disabling the LoRA adapter rather than loading a second model. $\lambda = 0$ is the unregularized objective and is the same code path the \texttt{fed} cells ran. The realized KL is reported beside the setting because a $\lambda$ that does not move the divergence is not a treatment. lambda=2 (0 of 2), lambda=5 (1 of 2), lambda=20 (1 of 2) not yet at depth and omitted.}
\begin{tabular}{@{}lrrrrr@{}}
\toprule
$\lambda$ & Cells & realized KL & mean $n_{\mathrm{ex}}$ & \texttt{self}$-$\texttt{fresh} & last 2 \\
\midrule
$0$ & 2 & -- & 39.1 & $+0.088$ & $+0.121$ \\
$0.05$ & 2 & 0.0047 & 42.1 & $+0.118$ & $+0.180$ \\
$0.2$ & 2 & 0.0033 & 40.8 & $+0.111$ & $+0.116$ \\
$1.0$ & 2 & 0.0023 & 23.6 & $+0.105$ & $+0.121$ \\
\bottomrule
\end{tabular}
\end{table}

\begin{table}[h]\centering\small
\label{tab:saturation}
\caption{A confident null, manufactured by a one-token change. \texttt{nb20} and \texttt{bk20} differ only in \texttt{KA\_KERNEL\_BANK}; model, seed, rounds, $k$, $n_{\mathrm{train}}=20$ and $n_{\mathrm{held}}=9$ are identical. Rounds are split on whether the frozen control has already saturated the scorer ($C_{\mathrm{reset}} \ge 0.49$; a correct kernel at parity scores exactly $0.50$), not on which bank they came from. Saturated rounds measure a mean absolute contrast of $0.0005$ against $0.1253$ where the instrument has room, a factor of $251$. The split holds \emph{within} \texttt{nb20}, and \texttt{nb20}'s unsaturated rounds agree with \texttt{bk20}'s, so the task bank does not change the effect --- it changes how often the instrument can resolve it. On its own, \texttt{nb20} reads as a tight equivalence result.}
\begin{tabular}{@{}llrrr@{}}
\toprule
Cell & Bank & saturated & mean $|\Delta|$ saturated & mean $|\Delta|$ has room \\
\midrule
\texttt{nb20} & default, 29 tasks & 8/16 & $0.0005$ & $0.1129$ \\
\texttt{bk20} & kernel, 124 tasks & 0/16 & --- & $0.1316$ \\
\midrule
pooled & --- & 8/32 & $0.0005$ & $0.1253$ \\
\bottomrule
\end{tabular}
\end{table}

\begin{table}[h]\centering\small
\label{tab:depth}
\caption{The same configuration read at two depths. \texttt{deep\_q15} and \texttt{fed\_q15} differ in one token of the command line, \texttt{--rounds 15} against \texttt{--rounds 5}; model, seed, $k$, $n_{\mathrm{train}}$, task bank and GPU map are identical. Every seed is positive over rounds 1--5 and negative over rounds 14--15, so stopping at five rounds and stopping at fifteen support opposite conclusions about whether self-training beats a frozen control. The loop is not starved at any point --- $n_{\mathrm{ex}}$ stays far above the starvation floor of two --- and round-0 retention falls to as little as $0.27$, which is the mechanism we attribute the reversal to (Section~\ref{sec:instrument}).}
\begin{tabular}{@{}lrrrrr@{}}
\toprule
Cell & Rounds & mean $\Delta$ r1--5 & mean $\Delta$ r14--15 & final retention & mean $n_{\mathrm{ex}}$ \\
\midrule
\texttt{q15\_s1} & 15 & $+0.180$ & $-0.183$ & 0.27 & 85 \\
\texttt{q15\_s2} & 15 & $+0.084$ & $-0.270$ & 0.26 & 69 \\
\texttt{q15\_s3} & 15 & $+0.082$ & $-0.089$ & 0.43 & 75 \\
\bottomrule
\end{tabular}
\end{table}

\begin{table}[h]\centering\small
\label{tab:lifetime}
\caption{When each arm reaches correct-at-parity, and what happens to the contrast afterwards. The \emph{gap} between the two arrival rounds bounds the largest contrast the design can ever observe, which is a property of the design rather than of the model. Where the arms arrive together the contrast never exceeds $0.062$ at any point in the run, however much data the learner receives: the \texttt{fedoc} cells take 350 verified examples per round and peak below that. Where they arrive two rounds apart it reaches $0.15$ to $0.27$. The \emph{final} column is near zero in every cell whose control has arrived, whatever the gap, because by the last round both arms sit at the ceiling -- so the gap caps how large an effect is visible, not merely when it stops being visible. Reported per cell and never pooled: the arrival rounds are integers on different models, and their mean would describe a cell that does not exist.}
\begin{tabular}{@{}lrrrrrrr@{}}
\toprule
Cell & Rounds & treat.\ at & ctrl.\ at & gap & peak & final & max $n_{\mathrm{ex}}$ \\
\midrule
\texttt{fed\_q05\_s1} & 5 & 2 & 4 & 2 & $-0.152$ & $+0.005$ & 275 \\
\texttt{fed\_q05\_s2} & 5 & 3 & -- & $>1$ & $+0.140$ & $+0.004$ & 257 \\
\texttt{fed\_q05\_s3} & 5 & 1 & -- & $>3$ & $+0.017$ & $+0.005$ & 304 \\
\texttt{fed\_q15\_s1} & 5 & 4 & -- & $>0$ & $+0.206$ & $+0.089$ & 104 \\
\texttt{fed\_q3\_s1} & 5 & 2 & 2 & 0 & $+0.042$ & $+0.014$ & 241 \\
\texttt{fed\_q3\_s2} & 5 & 2 & 4 & 2 & $+0.017$ & $+0.017$ & 194 \\
\texttt{fed\_q3\_s3} & 5 & 2 & 2 & 0 & $+0.029$ & $+0.006$ & 219 \\
\texttt{fedds\_s1} & 5 & 2 & 4 & 2 & $-0.272$ & $+0.002$ & 346 \\
\texttt{fedds\_s2} & 5 & 2 & 4 & 2 & $+0.183$ & $+0.011$ & 307 \\
\texttt{fedds\_s3} & 5 & 2 & 4 & 2 & $-0.136$ & $-0.002$ & 327 \\
\texttt{fedoc\_s1} & 5 & 1 & 1 & 0 & $+0.062$ & $-0.002$ & 350 \\
\texttt{fedoc\_s2} & 5 & 1 & 1 & 0 & $-0.006$ & $-0.002$ & 348 \\
\texttt{fedoc\_s3} & 5 & 1 & 1 & 0 & $-0.039$ & $-0.004$ & 348 \\
\bottomrule
\end{tabular}
\end{table}

\begin{table}[h]\centering\small
\label{tab:collapse}
\caption{Verified self-training past round five. The treatment arm peaks and then loses capability \emph{absolutely}: the drop column is its final score minus its peak. Two columns rule out the alternatives within each run. \texttt{trainC} is the model's score on its own training split and falls alongside the held-out score, which overfitting would not do. $C_{\mathrm{ctrl}}$ is the round0-replay arm --- identical LoRA, learning rate, schedule and round count, trained on verified self-data \emph{frozen at round 0} --- and it does not collapse, so the cause is the recursion rather than repeated training. It also reaches a higher score than the treatment ever does, so nothing is capping the treatment. Every training example in both arms passed execution against an fp32 reference.}
\begin{tabular}{@{}lrrrrrrrrrr@{}}
\toprule
Cell & Rounds & peak at & peak $C$ & final $C$ & drop & $C_{\mathrm{ctrl}}$ max & final & \texttt{trainC} & contrast @5 & contrast final \\
\midrule
\texttt{deep\_q15\_s1} & 15 & 4 & 0.47 & 0.18 & $-0.29$ & 0.48 & 0.48 & 0.50$\to$0.27 & $+0.182$ & $-0.251$ \\
\texttt{deep\_q15\_s2} & 15 & 4 & 0.49 & 0.16 & $-0.33$ & 0.49 & 0.49 & 0.49$\to$0.17 & $+0.071$ & $-0.321$ \\
\texttt{deep\_q15\_s3} & 15 & 2 & 0.48 & 0.24 & $-0.24$ & 0.51 & 0.49 & 0.49$\to$0.25 & $+0.082$ & $-0.097$ \\
\texttt{deepp\_q15\_s3} & 6 & 4 & 0.42 & 0.32 & $-0.09$ & 0.23 & 0.23 & 0.47$\to$0.47 & $+0.245$ & $+0.185$ \\
\bottomrule
\end{tabular}
\end{table}


\clearpage
\subsection{T1 --- Capability}
On the E0 gate (15 fixed Medium procedural tasks, $k{=}5$, 75 trials/model, 95\% Wilson~\cite{wilson1927interval} CIs), capability
separates into \emph{two walls}: correctness ranks the low/mid band, while \emph{speed}
(\texttt{fast\_rate}, the fraction beating the $\min(\text{eager},\texttt{torch.compile})$ roofline by
$\ge1.10\times$) is the frontier discriminator.

\begin{table}[H]\centering\small
\begin{tabular}{lccc}
\toprule
Model & correct & fast & mean $C$ \\
\midrule
GPT-5.6 Terra (api) & 0.920 & 0.667 & 0.793 \\
GPT-5.6 Sol (api)   & 0.907 & 0.600 & 0.753 \\
Kimi K2.5           & 0.387 & 0.240 & 0.313 \\
Fable 5.1           & 0.360 & 0.227 & 0.293 \\
Palmyra-X5          & 0.533 & 0.040 & 0.287 \\
gpt-oss-120b        & 0.467 & 0.080 & 0.273 \\
Qwen2.5-Coder-7B (open) & 0.280 & 0.000 & 0.140 \\
\bottomrule
\end{tabular}
\caption{T1 capability (top rows). \emph{Every} open model $\le$14B has $\texttt{fast\_rate}=0$: they can
sometimes write a correct kernel but essentially never a frontier-fast one; Coder-14B and Gemma-3-27B
additionally hallucinate \texttt{torch} internals. Source: \texttt{leaderboard.json}.}
\end{table}

\intuition{The two walls matter for everything above: T1 shows that below the frontier, models are
\emph{correctness}-limited long before they are \emph{speed}-limited. So the RSI rungs are mostly testing
whether a model can compound its way \emph{up the correctness wall}, not whether it can shave microseconds ---
which is why ``where $C$ moves, it moves via correctness acquisition, not speed'' recurs throughout.}

\subsection{T2 --- Weight-RSI: a rigorously-powered compounding null}
The title claim is where the strongest evidence lives, and it is a \emph{null}. Over multi-round rejection-sampling
SFT, carrying the lineage forward does not beat a matched reset.

\finding{The independent replicate is the \emph{trajectory} (one seed, one base checkpoint, one held
split, one accumulated adapter), not the round: rounds within a run are serially dependent, and an earlier
draft of this report pooled them as if independent. Pooled across scales (0.5B--14B) over
\textbf{57 trajectories} comprising 228 round-comparisons, \lmr{}\ $=-0.010\,[-0.035,+0.016]$, TOST-equivalent
at $\delta{=}0.05$ with $\bfz{=}5.6$ --- \emph{moderate} evidence for the null. A cluster-robust estimate
clustered on trajectory agrees ($-0.002\,[-0.024,+0.020]$, $\bfz{=}7.4$). The naive round-level pooling would
have reported $\bfz{=}14.7$; that figure is an artifact of the clustering and we do not claim it. The verdict
is robust: restricting to trajectories with $\geq3$ rounds ($n{=}36$) gives $\bfz\,{\approx}\,5.9$ and to
full-length $\geq6$-round trajectories ($n{=}29$) gives $\bfz\,{\approx}\,4.9$, both EQUIV.}

\begin{table}[H]\centering\small
\begin{tabular}{lccccc}
\toprule
Scale & traj.\ $n$ & rounds & \lmr{} (95\% CI) & TOST & \bfz{} \\
\midrule
Qwen2.5-Coder-0.5B & 12 & 51 & $-0.014\,[-0.058,+0.030]$ & EQUIV & 2.6 \\
Qwen2.5-Coder-1.5B & 20 & 93 & $-0.001\,[-0.037,+0.035]$ & EQUIV & 4.5 \\
Qwen2.5-Coder-3B   & 19 & 78 & $-0.027\,[-0.086,+0.032]$ & not equiv. & 2.7 \\
Qwen2.5-Coder-7B   & 5  & 5  & $+0.024\,[-0.126,+0.174]$ & underpowered & 2.0 \\
Qwen2.5-Coder-14B  & 1  & 1  & --- & n/a (single run) & --- \\
\midrule
\textbf{Pooled}    & \textbf{57} & \textbf{228} & $\mathbf{-0.010\,[-0.035,+0.016]}$ & \textbf{EQUIV} & \textbf{5.6} \\
\bottomrule
\end{tabular}
\caption{Lineage vs.\ matched reset at the \textbf{trajectory} level (primary estimand;
\texttt{equivalence\_tost.py}). $n$ counts independent trajectories; ``rounds'' counts the round-comparisons
they contain. Intra-trajectory ICC${=}0.093$ (design effect $1.28$), so the naive $n{=}228$ rounds carry the
evidence of ${\approx}178$ independent observations. Individual scales are now only moderately powered --- 3B
is \emph{not} individually equivalent once clustering is respected, and 7B ($n{=}5$) is not interpreted.
Per-round, \lmr{} shows a \emph{transient} early advantage (rounds 1--2) decaying to $\approx0$ by round 3.}
\end{table}

\noindent The uncontrolled T2 boards (\texttt{rsi\_leaderboard.json}, \texttt{tier\_speed\_rsi.json}) do show
apparent per-model ``compounding'' --- the strongest being \textbf{stable-code-3b}, self$-$fresh-frozen
$+0.226$, sign-consistent across 2 seeds; and starcoder2-15b $+0.074$ across 3 seeds. But the decisive
recursion-interruption fork (A1, \texttt{fork.json}) dissolves it: for deepseek-1.3b, continuing to self-train
vs.\ freezing the producer at the fork gives self$-$ckpt $=-0.001$ --- a \textbf{one-time upgrade, not
compounding}. Where $C$ does move it is almost entirely \emph{correctness acquisition}, not speed
(\texttt{decomp.json}: deepseek-1.3b $\Delta_{\text{correct}}+0.89$ vs $\Delta_{\text{speed}}+0.38$).

\intuition{The lineage-vs-reset null is the cleanest single fact in the project, and it is stronger than
``updates are small.'' If updates were merely small, lineage and reset would both drift slowly but lineage would
still inch ahead. Instead they are \emph{indistinguishable} --- which means the accumulated weight state carries
\emph{no useful cumulative information} beyond what one round of fresh self-training recovers. That is a
statement about the \emph{state}, not the step size.}

\subsubsection*{An H100 replication under a correctness-reliability metric}
The A100 compounding null above is measured on the headroom score. On H100 that score has no
range (Section~\ref{sec:instrument}), so we re-ran the same protocol under a pre-registered
second metric --- \emph{pass-rate}, the fraction of the $k$ sampled candidates that verify
(pre-registration, Amendment~1). Eight trajectories, 1.5B and 3B, two seeds per arm:

\begin{center}\small
\begin{tabular}{@{}lcc@{}}
\toprule
\textbf{Arm} & $\textsf{lineage}-\textsf{reset}$ & \textbf{95\% CI} \\
\midrule
control (published protocol) & $+0.184$ & $[+0.080, +0.287]$ \\
inject (teacher coverage on failed tasks) & $+0.184$ & $[+0.089, +0.278]$ \\
\midrule
inject $-$ control & $+0.000$ & --- \\
\bottomrule
\end{tabular}
\end{center}

\noindent Both CIs exclude zero: on this metric, accumulation beats a matched reset. This is the
first non-degenerate compounding measurement we obtained on H100, and we report it with three
limits that between them prevent it from supporting a compounding claim.

\textbf{First, the search baseline is degenerate under this metric and we do not report it.}
Best-of-$N$'s mechanism is the \emph{maximum} over $k(r{+}1)$ draws; pass-rate is a
\emph{mean}, which is invariant to the number of draws. Measured: as the budget grew
five-fold, $C_{\textsf{best-of-}N}$ moved $-0.024$. The apparent
$\textsf{lineage}-\textsf{best-of-}N = +0.676$ (positive in 37/37 rounds) would therefore
\emph{reverse} our search-beats-training result using an artifact. \textbf{A baseline is only a
baseline with respect to a scoring rule}; changing the rule silently changes what every
comparator means, and each must be re-derived rather than carried over.

\textbf{Second, the metric saturates at the top.} $C_{\textsf{lineage}} = 1.000$ in 18 of 37
rounds (49\%) on a five-task held set. Once lineage pins at the ceiling the difference measures
only how far \emph{reset} fell below it, so the magnitude is a lower bound and the
across-round trend is unreadable. Having diagnosed a floor effect and built a metric with range
at the floor, we landed in a ceiling effect: \textbf{saturation is a property of the match
between task difficulty and model ability, not of the metric}, and changing the metric only
moves where the wall sits.

\textbf{Third, the contrast is not training-compute matched.} Lineage retains its adapter and
trains every round ($R \times$ steps); reset re-initialises and trains on one round's data
($1 \times$ steps). Under a metric scoring reliability of correctness, ``more supervised
fine-tuning on verified-correct outputs raises the rate of correct outputs'' is close to
tautological. The result is evidence that accumulation improves correctness \emph{reliability};
it is not evidence of recursive self-improvement.

Finally, the positive control did not fail --- it ran out of work. Injected-task counts per
round run $[15,2,1,0,0]$ and $[19,3,0,0,0,0]$: injection targets tasks the student failed, and
by round~2 it solves everything, which is why the inject--control delta is exactly zero. A
five-task held set that a 1.5B model exhausts in two rounds cannot support a compounding claim
under any scoring rule. The remedy is the 456-task generated bank (Table~\ref{tab:reach}), not
a third metric.

\subsection{T3 --- Procedure-RSI}
When a frozen-weight model rewrites its own executable procedure, frontier models improve it substantially
(\texttt{trackc.json}): GPT-6 Astra $Q$ $0.298\!\to\!0.854$ (\textbf{+0.556}, growing over 6 rounds), Claude
Sonnet 5 \textbf{+0.401}, Mistral Large 3 \textbf{+0.310}; a frozen-procedure control (Fable-5.1) sits at
$+0.009$. But the self-modify mechanism (\texttt{selfplay\_mech.json}) shows the gain is
\textbf{68\% realized at round~0} (one-shot), \textbf{75\% non-recursive}, with the archive saturating by
round 1--2 and strategy sets already diverse (Jaccard 0.12).

\finding{The only models with \emph{sustained} per-round procedure gain (Astra, Mistral) are the ones that
started from \emph{low} $Q_0$ --- i.e.\ they are consuming headroom, not demonstrating recursive self-insight.
DeepSeek-V3.2's $-0.234$ ``self-degrade'' is a 1-round interrupted-run artifact and should not be read as a
finding.}

\paragraph{The one-shot reading is confounded by the harness, and we retract it.}
The procedure is a list of strategy strings, and the harness asks the model for ``$\le 12$'' of them while
displaying only 12 back. It is a \emph{soft} request --- the stored list is never truncated --- so what it
actually measures is instruction compliance. Frontier models comply exactly:

\begin{center}\small
\begin{tabular}{@{}llr@{}}
\toprule
\textbf{Model} & \textbf{Strategies per round} & \textbf{Total $\Delta Q$} \\
\midrule
GPT-6 Astra & 12, 12, 12, 12, 12, 12 & $+0.556$ \\
Claude Sonnet 5 & 12, 12, 12, 12, 12, 12 & $+0.401$ \\
Mistral Large 3 & 12, 12, 12, 12, 12, 12 & $+0.310$ \\
GPT-5.6 Sol & 12, 12, 11, 12, 11, 11 & --- \\
GPT-5.6 Terra & 10, 12, 12, 12, 12, 11 & --- \\
\bottomrule
\end{tabular}
\end{center}

\noindent Five of five interpretable runs sit at 11--12 from round~0; four of five never grow at all. A
procedure pinned at its limit in round~0 \emph{cannot be observed to keep improving}, so ``the model stopped
improving its procedure'' and ``the harness stopped letting it'' are not separable in this design, and the
one-shot interpretation is not earned by these data.

Re-running the same harness with an open model at caps $\{12, 48, \text{unbounded}\}$ shows procedures
\emph{do} grow when permitted --- 6 of 6 clean runs grow, two exceeding 12 strategies --- confirming that the
cap binds something real. It does \textbf{not}, however, produce a dose-response in quality: on the clean runs
mean $\Delta Q$ is $-0.046$ at cap~12, $+0.021$ at cap~48 and $-0.009$ unbounded ($n{=}2$ per cell, all six runs complete at six rounds). Lifting the cap lets the procedure grow without making it better.

We flag a correction here rather than silently reporting the clean figures. An earlier version of this
analysis reported a monotone rise to $\mathbf{+0.482}$ unbounded. That came from runs whose data directory sat
on an inode-exhausted filesystem, where the reference build raised \texttt{EDQUOT} and the grader returned
\texttt{False} for every candidate: those runs have $Q_0 = 0.000$ and four to five exactly-zero rounds, so
their apparent ``gain'' was recovery from a spurious zero baseline. We had already documented that failure
mode (it is defect~5 in Table~\ref{tab:defects}) and used the affected runs regardless. The exclusion is now
enforced in the analysis script rather than left to memory: any run with $Q_0 = 0$ or two or more
exactly-zero rounds is dropped, loudly.

We state the limits plainly. There are two seeds per cap cell; open-versus-closed is confounded with model
family; and because the cap is a soft request, the open model does not comply with it reliably, so these runs
corroborate rather than replicate. The strong evidence is the closed-model pinning, which is visible in the
\emph{already-published} run records and requires no new experiment. Settling the question properly requires
re-running frontier models at several caps. Until then we report the procedure gains as real and withdraw the
claim that they are one-shot.

\subsection{T4 --- Closed$\to$Open}
A closed researcher rewriting an open trainee's training harness can help substantially
(\texttt{combined\_rsi.json}): Fable-5.1 $\to$ Qwen-1.5B reaches improved$-$frozen \textbf{+0.472} (peak 0.472,
$C$ $0.147\!\to\!0.501$ vs frozen 0.029), and Fable-5.1 $\to$ Qwen-7B \textbf{+0.313}. But the sign is
\emph{pair-dependent}: the same researcher \emph{hurts} other trainees (Fable$\to$OpenCoder-1.5B $-0.265$,
Fable$\to$DeepSeek-1.3B $-0.444$), and many researcher rows are flat. As with T4's cousin T2, an open question
(ROADMAP A4) is whether the $+0.472$ is correctness rather than speed.

\intuition{Improvement \emph{does} transmit through authored tooling --- a closed model that cannot change its
own weights can still cause a better open model. But it is brittle and trainee-specific, which reads less like
``the researcher discovered a better training recipe'' and more like ``it happened to match this trainee's
correctness headroom.'' Transmitted improvement is real but not yet a reliable, general channel.}

\subsection{T5 --- Self-play: an author-yield failure, not a measured null}
The primary signal is \lf{} (extra gain from a co-evolving author beyond an escalating curriculum). \lf{} is
only defined if the live author actually produced accepted, valid, novel tasks. \textbf{It did not, at the
yield needed to measure anything.}

\textbf{Open} (\texttt{selfplay.json}, 24 runs): most models $\le$15B report $\lf\approx0$. Requiring at least
five accepted model-authored tasks before a row is interpreted, only 6/24 open runs qualify.
\textbf{Closed} (\texttt{selfplay\_closed.json}, procedure channel, 11 runs): 10/11 rows are \emph{exactly}
$\lf=0.0$ --- including runs whose authors proposed 36--37 tasks across 24 rounds. An exact zero across 24
rounds with non-zero proposal counts indicates the co-evolution channel never engaged, not that it engaged and
returned nothing. \textbf{Open-ended} (\texttt{open\_rsi.json}, 16 runs): all verdicts are ``one-time
(open $=$ fixed)''.

\finding{The two rows previously quoted as positives do not survive the yield floor: starcoder2-15b
$\lf=+0.232$ rests on \textbf{one} accepted authored task over 6 rounds, and deepseek-coder-6.7b-base
$\lf=+0.154$ on \textbf{three} over 3 rounds. Their per-round \lf{} traces oscillate in sign. We therefore
withdraw both from the headline claims and report T5 as \emph{undefined at the current author yield}.}

\noindent\textbf{What the dedicated diagnosis run actually shows.} The self-play diagnosis
(\texttt{selfplay\_diag.json}, Qwen2.5-Coder-3B, 8 rounds) authored \textbf{zero} valid tasks in every round.
Consequently every per-round diagnostic in that file --- authored solve-rate, trivial and unsolvable
fractions, authored headroom --- is \texttt{null}, and all three arms scored $C{=}0$. An earlier draft of this
report read that run as showing a \emph{bimodal} authored-difficulty distribution (tasks either trivial or
unsolvable, rarely in a learnable band). \textbf{That conclusion is not in the data and we retract it}: the
bimodality is a plausible hypothesis the run was unable to test, because nothing was authored to measure. The
supported statement is narrower: \emph{free-form task proposal collapses}. Without constrained, executable
task mutation, the author does not manufacture a frontier, so self-play cannot be tested at all --- which
makes structured mutation with validity and learnability gates a prerequisite experiment, not an extension.

\noindent (Note: an earlier README listed closed Sonnet-5 $+0.045$ / GPT-5.6-Sol $+0.041$ as ``emerging
positives''; these are \emph{not} in the current data and have been removed.)

% =====================================================================================
\section{Cross-cutting analyses: the mechanism}
% =====================================================================================
The rung-by-rung results tell us \emph{whether} improvement compounds. The following analyses tell us
\emph{why} it does not, and together they form the mechanistic core of the report.

\subsection{Coverage vs.\ scale, and the ``correctness wall'' as a sampling artifact}
Per-task $p$-maps on a fixed 29-task bank (\texttt{pmap\_curve.json}) trace coverage (tasks with $\ge1$
verified-correct sample) from 13.8\% at 0.5B to 86.2\% at 14B, and show that below the wall pass@$K$ dwarfs
pass@1:

\begin{table}[H]\centering\small
\begin{tabular}{lccccc}
\toprule
Size & $K$ & coverage & pass@1 & pass@$K$ & $K$-ratio \\
\midrule
0.5B & 64 & 13.8\% & 0.009 & 0.109 & ${\sim}12\times$ \\
1.5B & 64 & 27.6\% & 0.014 & 0.235 & ${\sim}17\times$ \\
3B   & 64 & 82.8\% & 0.112 & 0.777 & ${\sim}7\times$ \\
7B   & 48 & 82.8\% & 0.079 & 0.799 & ${\sim}10\times$ \\
14B  & 24 & 86.2\% & \textbf{0.655} & 0.862 & ${\sim}1\times$ \\
\bottomrule
\end{tabular}
\caption{Coverage and sharpening across scale. The sharpening geometry (results\_auto.tex): the
\emph{sharpenable} band ($0{<}p{<}0.2$) is 14\%/28\%/59\%/79\%/0\% at 0.5/1.5/3/7/14B --- it peaks mid-scale
and collapses at 14B, where all mass is already \emph{reliable} ($p\ge0.2$).}
\end{table}

\intuition{The sub-2B ``correctness wall'' is not an inability --- it is a \emph{low per-sample success
probability}. Give a small model enough draws and correct kernels appear (pass@$K \gg$ pass@1). This reframes
the entire problem: self-training does not need to \emph{learn to be correct}, it needs to \emph{reliably
reproduce} a correctness that is already latent in the sampling tail. That is a sharpening problem, not an
exploration problem.}

\subsection{Zero-score forensics: the wall is a formation failure}
Classifying \emph{why} each 0-score candidate failed (\texttt{forensics\_summary.json}) shows the sub-3B wall is
a \textbf{kernel-formation} wall (no parseable \texttt{ModelNew}), not runnable-but-wrong code: \texttt{no\_extract}
dominates at 0.5B (72\%), DeepSeek-1.3B (90\%), Yi-1.5B (96\%). The failure \emph{locus marches downstream with
scale}: incoherence $\to$ truncation/syntax (Qwen-7B: 75\% syntax) $\to$ API-hallucination $\to$ wrong-output
$\to$ correct (Qwen-14B: 65\% correct, only 0.2\% \texttt{no\_extract}).

\intuition{Small models do not fail by writing subtly-wrong kernels --- they fail by not producing a
well-formed kernel at all. That is why coverage is so low and why self-training starves: there is nothing
correct to train on. As scale rises the bottleneck migrates from ``can it write a kernel'' to ``is the kernel
right'' to ``is it fast.''}

\subsection{Search beats training at matched compute}
Analysed at the trajectory level over 57 trajectories / 228 round-comparisons
(\texttt{search\_vs\_train.json}), lineage$-$bestof$N = -0.111\,[-0.138,-0.083]$ (CI excludes zero; the
cluster-robust estimator, which weights rounds rather than runs, gives $-0.142\,[-0.163,-0.120]$ --- a larger effect, same sign, and its interval also excludes zero), and verified search beats lineage in \textbf{84\% of runs}
(76\% of rounds); every well-powered scale is individually negative (0.5B $-0.123$, 1.5B $-0.114$,
3B $-0.123$; 7B $-0.018$ at $n{=}5$ is not interpreted). Unlike the compounding null, this positive
\emph{strengthens} under the clustering correction, because the contrast is made within a round. A separate matched-budget baseline sweep (\texttt{baselines.json},
budget 20) is weaker and mixed --- Qwen-1.5B recursion\_gain $+0.055$ over its best non-recursive baseline, but
starcoder2-3b $-0.073$ --- so the strong, consistent statement is the search-vs-lineage one.

\finding{At a matched compute ledger, best-of-$N$ verified search beats lineage self-training (lineage$-$bestof$N
= -0.111\,[-0.138,-0.083]$ over 57 trajectories, CI excludes zero). The useful solutions are \emph{already reachable by sampling};
search exploits them directly, while training must compress them into shared weights and have the gain survive.}

\subsection{WHY-RSI: the weight-level mechanism and the inverted-U in scale}
Across 133 weight-RSI runs (\texttt{mech\_analysis.json}) the correctness wall is the primary gate: crossing
rises from 54\% below 2B to 92\% at $\ge$2B. Among the 97 wall-crossers, RSI is an \textbf{inverted-U in scale}:

\begin{table}[H]\centering\small
\begin{tabular}{lccccc}
\toprule
Band & $n$ & wall-cross & LoRA drift & retention & frac RSI \\
\midrule
$<$2B  & 67 & 0.54 & 0.019 & 0.013 & 0.03 \\
2--8B  & 40 & 1.00 & 0.271 & 0.180 & \textbf{0.50} \\
$\ge$8B & 26 & 0.81 & \textbf{0.596} & 0.339 & 0.23 \\
\bottomrule
\end{tabular}
\caption{The inverted-U (\texttt{mech\_analysis.json}, $n{=}133$). RSI peaks at 2--8B; the largest models drift
\emph{most} in weight space yet compound \emph{least} --- motion without progress near the roofline. 28/133
runs met the compounding criterion (stable-code-3b recurs). When drift occurs it localizes to \emph{late} layers (task
adaptation) and the held-family transfer gap is $\approx0$ (no transferable meta-skill).}
\end{table}

\finding{A subtle and important nuance: on this bank, compounders have \emph{lower} generation diversity than
flat runs (0.41 vs 0.52), not higher. The discriminator between compounding and flat among wall-crossers is
\textbf{sustained LoRA drift + retention} (${\sim}5\times$ higher in compounders), \emph{not} diversity
preservation. So ``diversity collapse causes the null'' is \emph{not} supported by our weight-RSI data ---
compounding here looks like \emph{productive convergence onto a good basin}. Diversity contraction remains a
plausible \emph{cross-round option-value} lock (see the Discussion) but we flag it as an association whose clean
test is a diversity-preserving intervention, and we explicitly do \emph{not} claim it as the mechanism.}

\intuition{Scale does not monotonically help RSI --- it is an inverted-U. Small models cannot form correct
kernels often enough to feed the training set (formation-starved). Mid-scale models sit on the widest
\emph{learnable frontier} and compound best. Large models cover almost everything already, so their weights
churn (high LoRA drift) without moving the held-out score --- ``motion without progress'' near the roofline.
The binding constraint slides from \emph{formation} (small) to \emph{headroom} (large).}

\subsection{Probe-as-intervention: correctness is represented better than it is generated}
A linear probe reads ``is this kernel correct'' out of mid-to-late layers at high AUC (Qwen-3B best\_auc
\textbf{0.985} at depth 0.86, while its actual correctness rate is 0.105) --- the model \emph{represents}
correctness far better than it \emph{generates} it (\texttt{interp.json}).

\finding{Analysed at the right unit, probe-guided selection \emph{does} beat uniform-random selection at
matched \emph{generation} budget. Clustering the 19 usable runs on their 12 distinct base checkpoints (rows
sharing a checkpoint are repeated draws, not distinct models), the mean lift is
$\mathbf{+0.124\,[+0.029,+0.218]}$ --- the 95\% CI excludes zero --- recovering $36\%\,[11\%,60\%]$ of the
oracle$-$random headroom. An earlier draft read the same rows as ``no advantage'' because it looked at the
spread of individual runs ($-0.27$ to $+0.45$) rather than the clustered mean. See
Appendix~\ref{app:probe}.}

\noindent\textbf{Three reasons this is a modest secondary result, not a headline.} (i) \emph{Uniform random is
the weakest possible comparator}: the informative ones --- length-normalised log-likelihood reranking, a
compile/run validity filter, execution agreement --- were not run, so we cannot say the probe is the best use
of the same compute. (ii) \emph{The denominators are tiny}: each run's lift is computed over 5--12 test tasks,
so a ``$+0.45$'' is about five problems, and the clustered interval is correspondingly wide. (iii)
\emph{Matched-generation selection is not matched-verification harvesting}: a gain in top-1 correct rate from
reranking an existing pool does not imply the probe makes a verification-constrained search cheaper.

\noindent\textbf{What we cannot verify from the released artifacts.} Two figures quoted in earlier drafts --- a
within-task ranking AUC of ${\approx}0.67$ and a matched-verification harvesting yield of ${+}0.01$ --- come
from a probe-baselines run whose per-candidate outputs are \emph{not} in this repository. Per-candidate probe
scores were not retained anywhere, so the within-task AUC cannot be recomputed from what we release. We
therefore do not rely on either number, and we state the within-task estimand as \textbf{unmeasured} rather
than as measured-and-small. This is the single most important gap in the probe story: only a within-task
ranking metric can explain a selection gain mechanistically, and the pooled AUC cannot substitute for it ---
as the runs with pooled $\mathrm{AUC}{=}1.000$ and lifts anywhere from $+0.000$ to $+0.250$ demonstrate.

\intuition{A linear probe can read ``is this kernel correct'' out of mid-to-late layers at high AUC, yet the
model emits a correct kernel far less often than it can internally recognize one. The model \emph{knows} more
than it \emph{generates}. This is the same discovery$\to$assimilation gap seen from the inside: the limiting
step is turning latent competence into a reliably-produced output.}

% =====================================================================================
% =====================================================================================


\clearpage
% The heading and label live here, not in the included file: the old root supplied them, and
% without them this section printed untitled and two \ref{sec:instrument} in the synthesis
% pointed at nothing.
\section{Instrument validity: how a self-improvement benchmark manufactures a null}
\label{sec:instrument}
% =====================================================================================
% Instrument validity. Included from kernelascent_full.tex before the Synthesis, because it
% qualifies every null the Results section reports.
% =====================================================================================
A benchmark that reports a null owes its reader a second result: evidence that the
measurement \emph{could have} reported a positive. Self-improvement benchmarks are unusually
exposed here, because the headline quantity is a difference between two arms of the same
pipeline. Any defect that suppresses both arms equally leaves the contrast at zero while every
surface diagnostic --- jobs complete, kernels compile, scores populate --- looks healthy. A
broken instrument does not report that it is broken; it reports $\Delta \approx 0$.

Over this project, \textbf{five independent pipeline defects each produced a plausible,
publishable-looking null}, summarised in Table~\ref{tab:defects}. None announced itself, and
three survived a code review before being caught by an assertion.

\begin{table}[t]
\centering\small
\begin{tabular}{@{}p{0.30\linewidth}p{0.31\linewidth}p{0.31\linewidth}@{}}
\toprule
\textbf{Defect} & \textbf{Presented as} & \textbf{Caught by} \\
\midrule
Grader mis-imported under a wrong root & ``self-training does not compound'' & identity-kernel assertion \\
Grading device index absent on a 1--2 GPU allocation & ``the model writes bad kernels'' & same assertion, run \emph{on the allocation shape the jobs use} \\
Raw generations graded without extracting the submission & ``the 14B teacher solves 0/29'' & non-zero coverage precondition \\
Score saturated at the correctness floor & ``lineage $\approx$ reset'' & score-histogram check \\
Compiled-baseline cache hitting a filesystem inode quota & ``4 of 6 rounds produced nothing'' & output-path/quota correlation \\
\bottomrule
\end{tabular}
\caption{Five defects, each of which independently manufactured a null. The middle column is
what the artifact would have claimed had the defect not been found.}
\label{tab:defects}
\end{table}

\paragraph{A clean zero is a bug hypothesis before it is a finding.}
The operational signature that caught most of these is distributional, not statistical. A
\emph{real} null is noisy: it has per-task spread, partial successes, seed-to-seed variance. A
null in which every task fails identically, or every score takes the same value, is
instrumentation. We therefore treat suspicious tidiness as evidence \emph{against} the result,
inverting the usual intuition that clean data is good data.

\paragraph{Correctness of the instrument is not sufficient; it must also have range.}
This is the failure mode we consider most transferable, because it defeats the obvious check.
Our headroom-normalised score is
$0.5 + 0.5\cdot\mathrm{clamp}\!\left(\frac{sp-1}{\mathrm{ceiling}-1}\right)$,
which is well-formed, was verified end-to-end on known-good inputs, and was working exactly as
specified. It was nonetheless dead as a measurement: on H100, \textbf{76\% of all scores across
every run and round were exactly $0.50$} --- precisely the value of a correct kernel at
parity with the baseline. Models reach the correctness floor in one to three rounds and stop,
because essentially no model from 0.5B to 14B produces a kernel faster than \texttt{torch.compile}
on this hardware.\footnote{The 14B teacher's best speedups across the bank: nine tasks at
$1.00\times$, two at $1.08\times$, and one each at $1.05\times$, $1.01\times$, $2.03\times$.}
With both arms pinned to the same absorbing value, $\textsf{lineage}-\textsf{reset}$ is forced
to zero \emph{by construction}, and the experiment could not have detected compounding in
either direction.

Critically, this is not repairable by task curation. Filtering the bank against an
H100 anchor with a $1.3\times$ minimum roofline headroom kept 23 of 29 tasks and still returned
a frozen-base mean best score of $0.511$, with every surviving task at $\approx 0.50$ and every
L3 task at $\textrm{correct\_rate}=0$. The difficulty distribution is bimodal with no middle
band: tasks are either unreachable or solved-at-parity. A filter selects among existing tasks;
it cannot create a regime in which the model is correct \emph{and} has headroom it can capture.

\paragraph{The estimator, not the normalisation, is the thing to fix.}
A saturated metric invites the wrong repair. The instinct is to re-normalise; the actual
problem is that best-of-$k$ over a speed score is the wrong \emph{estimator} for a population
sitting at the correctness floor. Best-of-$k$ answers ``can it ever,'' and returns an identical
$0.500$ whether a model solves a task once in eight attempts or eight times out of eight ---
two plainly different models. Where the live axis is reliability rather than speed, pass-rate
over $k$ separates those cases by a factor of eight in resolution.

We report such runs under a separate metric (pre-registration, Amendment~1), never pooled with the
headroom leaderboards, and with two properties stated up front: the denominator is $k$, so a
model emitting nothing parseable scores zero rather than being excluded; and it does not reward
speed, so it tests a \emph{weaker and different} claim than the headroom score. It was also
adopted \emph{after} observing saturation --- a forking path that we record as such rather than
present as the original design.

\paragraph{Check the denominator of the mechanism, not just the metric.}
The same failure shape recurs wherever a rate is computed over events that never occurred. Two
frontier models in our procedure-RSI track recorded $n_{\textrm{strategies}}=0$ in every round:
their output was never parsed into a strategy list, so the self-modification channel never
engaged. The resulting $-0.234$ ``self-degradation'' was widely quoted from an earlier draft of
this work and is a parse failure. Likewise, our self-play primary metric requires the live
author to produce accepted valid tasks; demanding at least five qualifies only 6 of 24 open and
5 of 11 closed runs. \textbf{A metric computed over zero events is undefined, not zero}, and the
distinction is invisible in the aggregate.

\paragraph{What we propose as standard practice.}
We ship these as executable gates rather than advice, and we recommend them for any
self-improvement benchmark reporting a null:
\begin{enumerate}[leftmargin=1.6em,itemsep=1pt]
\item \textbf{Known-good input, known-good output}, asserted before every batch and run
\emph{on the same allocation shape as the experiment} --- one of our defects lived entirely in
the difference between a 1-GPU and a 3-GPU allocation.
\item \textbf{A dynamic-range precondition}: demonstrate the primary metric takes at least two
distinguishable values on this hardware and this task bank, \emph{before} interpreting a
contrast computed from it.
\item \textbf{A positive control on the causal path}: an intervention known to move the primary
metric (for us, injecting verified solutions on tasks the student failed). A null is
interpretable only once the harness has been shown to register a positive.
\item \textbf{Mechanism denominators reported alongside every rate}, so undefined is
distinguishable from zero.
\item \textbf{Automated cross-artifact consistency checking}: every number in the manuscript
re-derived from the data files at build time. Ours catches stale figures and claim strings that
contradict the data they cite --- both of which occurred.
\end{enumerate}

We emphasise that this is not a catalogue of unusual carelessness. Each defect was a single
line, each was locally reasonable, and each was caught only because an assertion was looking
for it. The generalisable claim is structural: \textbf{in a benchmark whose headline is a
difference between two arms of one pipeline, the cheapest result to produce accidentally is a
null}, and the field currently has no convention requiring authors to rule that out.


% discussion.tex is NOT input here: 05_synthesis.tex already inputs it, demoted to a
% subsection under the Synthesis header. Inputting it in both places typeset the whole
% mechanism analysis twice, four pages apart, under two identically-titled headings.
\clearpage
\section{Synthesis: why self-improvement does not compound here}
% =====================================================================================
% discussion.tex carries its own \section + \paragraph structure; demote its \section to a
% subsection locally so it nests cleanly under this Synthesis header.
{\let\section\subsection % Discussion: why verified self-training does not compound (distilled from GPT-6-Astra analysis, 2026-09-18).
% Static, hand-checked. \input by figures.tex.
\section{Why self-improvement does not compound}

\paragraph{The mechanism: a discovery$\to$assimilation bottleneck (sharpening, not exploration).}
Our two strongest results jointly localise the failure. Best-of-$N$ search beating lineage training at matched
compute (lineage$-$bestof$N=-0.111\,[-0.138,-0.083]$ over 57 independent trajectories), together with the
large pass@$K\!\gg\!$pass@1 gap, shows
that \emph{useful solutions are already present in the sampling tail} but are not reliably produced in one
sample. The limiting pipeline is therefore
\[
\text{rare verified sample}\;\to\;\text{selected SFT example}\;\to\;\text{weight update}\;\to\;\text{reliably reusable gain},
\]
and it is \emph{lossy}: search can exploit a rare discovery directly, whereas training must compress it into a
shared parameterisation, preserve other skills, and have the gain survive later rounds. The lineage-vs-reset
null (inheriting updates does not beat restarting) is direct evidence \emph{against a useful cumulative learning
state}, not merely evidence that individual updates are small. This is the classic self-training signature:
selected-sample training concentrates probability on \emph{already-reachable} successes faster than it expands
what is reachable (cf.\ STaR, ReST, ReST$^{EM}$, expert iteration; and the RLVR capability-boundary debate,
Yue et al.\ 2025).

\paragraph{A formal fixed point --- with its limits stated precisely.}
Under an \emph{idealised} support-preserving rejection-conditioning operator $R$ (unlimited accepted samples,
exact per-task fitting, no cross-task transfer), exact solution-support coverage
$C_0(\pi)=\Pr_x[p_\pi(x)>0]$ is invariant and $R$ is idempotent after one step: $C_0(R\pi)=C_0(\pi)$ and
$R^2\pi=R\pi$. On tasks with positive success probability $R$ places all mass on accepted outputs; on $p{=}0$
tasks no accepted distribution exists. So an idealised verified-SFT loop is a \emph{one-step fixed point in
exact coverage}. \textbf{We are careful about what this does \emph{not} say.} (i) Neural SFT is not this
operator --- shared parameters can generalise and raise the probability of unsampled outputs, so real SFT need
not preserve support. (ii) \emph{Finite-budget} coverage is \emph{not} invariant:
$C_K(\pi)=\mathbb{E}_x[1-(1-p_\pi(x))^K]$ rises sharply under sharpening even when $C_0$ is fixed --- ``pass@$K$
coverage is a fixed point'' is false. (iii) Zero observed successes only bound $p\lesssim 3/K$; our
``unreachable mass'' is more precisely \emph{undiscovered mass at the evaluated budget}. (iv) A training-operator
fixed point is \emph{not} a parameter-space local optimum. We therefore report an \textbf{operational
self-training plateau / protocol-relative fixed point}, never an unqualified ``local optimum.''

\paragraph{Why verified search escapes where training does not.}
Search does not escape the policy's support; it escapes the \emph{one-sample} bottleneck and avoids the
assimilation bottleneck. For success probability $p$, $P(\text{hit in }N)=1-(1-p)^N$ makes extra draws
valuable when $p$ is small, and the verifier keeps the winner. Search (a) uses a discovery directly rather than
compressing it into weights, (b) allocates candidates per-task rather than via one shared update, and (c) does
not permanently contract the proposal policy. The precise statement: \emph{verified search accesses useful
existing tail capability more efficiently than the tested training loop consolidates it}. (This does not settle
lifetime economics --- a training update could amortise over enough future queries; that needs a separate
deployment-horizon study.)

\paragraph{Diversity contraction as a reinforcing lock (associational).}
Selected-sample SFT can drive a loop: current common winners $\to$ accepted set $\to$ higher probability on
those winners $\to$ fewer alternatives in the next accepted set. The loss is \emph{option value} --- rare
decompositions and kernel families that would enable \emph{later} discoveries. If undiscovered useful mass
$U_t$ shrinks, later rounds find fewer novel successes ($1-(1-U_t)^N\le NU_t$) even as pass@1 rises on familiar
ones. Higher retention among compounders and diversity collapse at large scale are \emph{consistent} with this
lock, but we flag them as association, not proof: low diversity can also reflect successful specialisation, and
the cleanest test is a matched-reward, matched-example, diversity-\emph{preserving} intervention.

\paragraph{The inverted-U, mechanistically.} Partition tasks into \emph{undiscovered} ($p{\approx}0$ at budget),
\emph{learnable frontier} ($0<p<0.2$), and \emph{near-ceiling} ($p$ high). Most self-training value comes from
the middle band. Small models: formation failures ($P(A\mid x)=P(F\mid x)P(A\mid F,x)$ with tiny $P(F)$) starve
the accepted set; mid-scale: the frontier is widest, so discoveries become signal (RSI peaks); large scale:
high coverage but little remaining headroom, so drift redistributes mass within already-covered behaviour
(``churn without gain'' --- supported descriptively; ``only superficial late-layer edits'' is \emph{not} yet
established, as parameter distance is not calibrated functional change).

\paragraph{What would break the plateau, and what our data predict.}
Two distinct breakouts are needed: \emph{accumulation} (retain/express more already-discoverable successes ---
could yield lineage$>$reset) and \emph{frontier expansion} (make new strategies operationally reachable ---
required once the frontier is exhausted). A productive loop needs
$\text{new discoveries}+\text{assimilation}+\text{retention}+\text{headroom}$ simultaneously. The single most
informative next experiment is a \textbf{coverage-injection $\times$ diversity-preserving} factorial: inject
stronger-model / repair / search-generated solutions on \emph{undiscovered} tasks and test whether inherited
gains then emerge --- separating ``lack of discoveries'' from ``loss of discoveries after learning.''
Exploration bonuses, novelty-selected accepted sets, prerequisite curricula, off-policy replay, and
search-amplified expert iteration are ranked candidates; each helps only the specific deficit it targets (e.g.\
entropy cannot fix formation; off-policy replay of redundant winners adds no frontier information).

\paragraph{Scope --- claims we do \emph{not} make.} We do not claim the models' reachable set is mathematically
invariant, that the weights sit at a genuine local optimum, or that RL/RLVR cannot expand capability. We claim,
with \emph{moderate} evidence for the null and a significant positive respectively: under verified
rejection-sampling SFT at the tested scales/budget, carrying the lineage forward does not beat a matched reset
(trajectory-level TOST equivalence at $\delta{=}0.05$, $\mathrm{BF}_{01}{=}5.6$ over 57 independent
trajectories); verified search is a better use of the same compute (CI excludes zero); and the binding internal signals are coverage/formation (small), a mid-scale learnable
frontier, and drift-without-gain plus diversity contraction (large). The synthesis: \emph{expert iteration
without sufficient expert amplification, plus incomplete retention of the apprentice's useful tail.}
}

% =====================================================================================
\section{The accumulated intuition (what three weeks taught us)}
% =====================================================================================
This section is the part that does not fit in a conference paper: the qualitative model we now carry in our
heads about RSI in this domain.

\begin{enumerate}[leftmargin=1.6em]
\item \textbf{Self-training is a sharpening operator, not an exploration operator.} It concentrates probability
on already-reachable successes faster than it expands what is reachable. Everything else follows from this.

\item \textbf{Coverage is the currency, and it is not self-generated.} A round of self-training can only reuse
solutions the model already finds; it cannot manufacture coverage on tasks where the per-sample success
probability is effectively zero at the evaluated budget. Compounding is coverage-limited, and coverage is
exogenous to the loop.

\item \textbf{The bottleneck is discovery$\to$assimilation, not knowledge.} The model represents correctness it
cannot reliably generate (probe result), and search --- which uses discoveries directly --- beats training,
which must assimilate them into weights and retain them. The lossy step is assimilation-and-retention.

\item \textbf{``Operational plateau,'' not ``local optimum.''} We are careful: the null is a protocol-relative
fixed point of the tested verified-SFT loop, not a claim that the weights sit at a parameter-space local
optimum. Finite-budget coverage is \emph{not} invariant; a better loop could still move.

\item \textbf{Diversity contraction is a plausible reinforcing lock.} Selected-sample SFT can raise probability
on current winners at the cost of the rare decompositions that would enable \emph{later} discoveries. We flag
this as association, not proof --- the clean test is a diversity-preserving intervention.

\item \textbf{The open contribution is a better \emph{harness}, not a better model.} Because RSI is a protocol
we run, the way to break the plateau is to change the protocol: inject coverage on undiscovered tasks, preserve
diversity in the accepted set, amplify the ``expert'' with search before assimilation. This is exactly what the
public benchmark's harness-submission track exists to surface.
\end{enumerate}

% =====================================================================================
\section{What would break the plateau: predictions}
% =====================================================================================
Two distinct breakouts are needed and they target different deficits: \emph{accumulation} (retain and express
more already-discoverable successes --- could make lineage finally beat reset) and \emph{frontier expansion}
(make genuinely new strategies operationally reachable --- required once the current frontier is exhausted). A
productive loop needs new discoveries $+$ assimilation $+$ retention $+$ headroom \emph{simultaneously}. Our
single most informative next experiment was, at the time of writing, a \textbf{coverage-injection $\times$
diversity-preserving} factorial. The instrument work in Section~\ref{sec:instrument} has since superseded that.
The coverage-injection control was run, and its injected-task count falls to zero by round~2 because the student
solves the entire held set; injection cannot separate ``lack of discoveries'' from ``loss of discoveries'' on a
task the model has already saturated.

The informative experiment is now the one a repaired grader makes possible for the first time: run the
compounding protocol on \textbf{kernel authoring} rather than correctness-preserving rewriting. Under the
published prompt, models attempt a custom kernel in \textbf{0 of 16} samples; when asked explicitly, \textbf{14
of 16}, of which \textbf{1} verifies. A success rate in that band is far from both the floor and the ceiling
that defeated every metric we tried, which is precisely the regime a compounding study requires. Whether
self-training compounds on a task the model has \emph{not} saturated is, as far as we can tell, still open ---
and it is a question this benchmark could not previously have asked.

% =====================================================================================
\section{Limitations and scope}
% =====================================================================================
We claim, with moderate evidence on the conventional thresholds~\cite{kass1995bayesfactors} ($\mathrm{BF}_{01}=5.6$ at the trajectory level): under verified rejection-sampling SFT~\cite{zelikman2022star,gulcehre2023rest} at the tested scales and budgets, carrying
the lineage forward does not beat a matched reset; verified search is a better use of the same compute; and the
binding internal signals are coverage/formation (small scale), a mid-scale learnable frontier, and
drift-without-gain plus diversity contraction (large scale). We do \emph{not} claim the reachable set is
mathematically invariant, that the weights sit at a genuine local optimum, or that RL/RLVR (as opposed to
rejection-sampling SFT) cannot expand capability. All primary intervals are now computed at the
\emph{trajectory} level rather than the round level, since rounds within a run share a base checkpoint, a seed,
a held-out split and an accumulated adapter; this is the honest unit of replication and it costs us evidential
strength, moving the equivalence Bayes factor from $\mathrm{BF}_{01}=14.7$ (pooling 228 rounds) to $5.6$
(57 trajectories) --- strong to \emph{moderate}. The verdicts survive the correction, and the
search-beats-training positive strengthens under it, because that contrast is made within a round. Per-trajectory
(2--3 seed) uncertainty remains the binding constraint. Finally, all results here are A100; our H100 runs are
reported separately and are non-comparable on two axes at once: different roofline constants, and a harness
that silently rejected every non-PyTorch submission (Section~\ref{sec:instrument}). An earlier version of this
sentence blamed a stronger \texttt{torch.compile} baseline; that attribution was wrong --- the default scorer
never used the compiled baseline --- and is retracted in Section~\ref{sec:instrument}.

% =====================================================================================
\section{Conclusion}
% =====================================================================================
Three weeks of clean-verifier, compute-matched experiments across five rungs of recursion and scales from 0.5B
to 14B converge on one picture: in GPU-kernel optimization, verified self-training \emph{sharpens what a model
already covers but does not explore}, so improvement does not compound --- not because the models lack
capability, but because the loop's feedback channel assimilates rare discoveries lossily and cannot manufacture
the coverage it would need to keep going. The result is a strong, honest null flanked by two positive findings
(search beats training; the correctness wall is a sampling artifact) that together localize the failure and
point at the intervention most likely to break it. KernelAscent ships that intervention as an open, automated
submission track, so the community can try to turn the null into a curve.


\clearpage
\bibliographystyle{plain}
\bibliography{refs}

\clearpage
\appendix
\section{What was removed, narrowed, or retracted --- and why}
\label{sec:removed}

Every claim this project has withdrawn or restricted is listed here with the evidence that forced
the change. Nothing is dropped silently. A report that shows only its surviving results has
removed exactly the information a reader needs in order to calibrate the rest.

\subsection{Retracted: claims the data does not support}

\begin{description}[leftmargin=1.4em,style=nextline]
\item[{T1-kernel monotone ordering ($p = 1/120$).}]
  The perfectly ordered five-point decline was an artifact of \texttt{KA\_EXTRACT=strict}, a
  candidate filter that discards 66\% of 0.5B generations and 5\% of 14B generations --- a
  severity that correlates with the independent variable. Under the permissive policy Spearman
  $\rho$ moves from $-1.000$ to $-0.900$ and 3B breaks the order. \textbf{The endpoint contrast
  survives both policies and is retained}; the ordering claim and its $p$-value are withdrawn.

\item[{T1-kernel ``14B verifies nothing''.}]
  Recorded twice from a cell that was still running (0 verified at 108 and at 252 attempts). The
  finished cell has 2 verified in 338. 14B is near-zero, not zero, and the falsifier language now
  says so.

\item[{T3 ``frontier self-modification is one-shot''.}]
  The harness requests at most 12 strategies; 5 of 5 interpretable runs sit at 11--12 from round
  0 and 4 of 5 never grow. The observed plateau is the cap, not the model. Gains against a frozen
  procedure remain real and are reported.

\item[{T3 cap dose-response ($+0.482$ at unbounded).}]
  Came from runs carrying the EDQUOT signature ($Q_0 = 0.000$ and four to five all-zero rounds), a
  filesystem failure in which the grader raises and every candidate scores zero. Clean runs show
  no dose-response: $-0.046 / +0.021 / -0.009$.

\item[{``H100's \texttt{torch.compile} closed the headroom''.}]
  Offered as the explanation for a spike of scores at exactly $0.50$. It is false: the default
  scorer never touched the compiled baseline at the time. The spike is a plain-torch rewrite
  landing at parity, which is a different and more consequential fact.

\item[{Teacher injection improves compounding.}]
  First reported as $+0.100$ against $+0.051$ for the control. It does not survive dropping
  rounds whose lineage was severed by a resume defect, nor matching on round index, since the
  arms had unequal valid depth. Matched: control $+0.057\,[+0.020,+0.094]$ versus inject
  $+0.075\,[+0.035,+0.114]$; and on the completed pass-rate board, $+0.440$ versus $+0.445$.
  \textbf{Not detected.}

\item[{``Retrieval is the worst baseline''.}]
  Stated from one cell at $0.154$. The completed four-cell sweep puts retrieval at $0.154$,
  $0.329$, $0.385$ and $0.426$ --- a spread of $0.272$, against within-cell gaps between any two
  methods of at most $0.27$ and typically half that. It is the worst method in one cell and
  within $0.013$ of the best in another. Retrieval is high variance, not reliably weak, and the
  original claim was an artifact of reading the low end of that range as the value.

\item[{``76\% of all scores were exactly 0.50''.}]
  Not reproducible from any artifact. The defensible figure is that 75\% of \emph{non-zero}
  scores fall within $\pm 0.01$ of parity on the 29-task calibration set.
\end{description}

\subsection{Narrowed: claims that hold with a smaller scope}

\begin{description}[leftmargin=1.4em,style=nextline]
\item[{Held-family transfer, on the DSL kernel bank only.}]
  The transfer measure groups tasks by the prefix before the first underscore, which separates
  the legacy bank into real op-families ($\texttt{l2\_7}\rightarrow\texttt{l2}$, three or more
  families) and collapses entirely on the DSL kernel bank, where every task is named
  $\texttt{dsl\_}\langle\text{op}\rangle\texttt{\_}\langle\text{dtype}\rangle\texttt{\_}\langle\text{shape}\rangle$.
  There the whole 124-task bank is a single family, so the ``held-out family'' is 20 arbitrary
  tasks drawn from the training distribution and $C_{\mathrm{transfer}}$ measures ordinary
  held-out generalisation while carrying a label that claims more. A transfer gap of
  $\approx 0$ is then the expected outcome and is not evidence about a transferable meta-skill.
  The $n=133$ mechanism analysis that states that gap runs on the legacy bank, whose split is
  not degenerate, so that claim stands as written; no kernel-bank cell contributes a transfer
  figure to this report. The lab now records \texttt{degenerate\_family\_split} in every
  artifact so the condition travels with the data. We did not change the grouping rule while
  cells were in flight: these runs resume from checkpoints, and a cell that resumed under a new
  split would have early and late rounds measured on different sets, which is a worse failure
  than the one being fixed.

\item[{The compounding null ($n=57$ trajectories, \bfz{}$\,{=}\,5.6$).}]
  Valid, and narrowed twice. It is a null about acquiring \emph{correctness} rather than
  optimisation, and specifically about \textbf{best-of-$k$} correctness, which by construction
  cannot see a reliability gain. The pass-rate board finds a large, family-transferring effect on
  that axis. These are not contradictory: the null is valid and narrow, and its narrowness was
  invisible until a metric existed that could see past it.

\item[{T1-kernel shape across scale.}]
  Reported as a decline with 14B as the endpoint. The completed 32B sweep recovers to $2.99\%$
  from 14B's $0.00\%$ (Fisher $p = 0.015$), so the curve is a \textbf{U} with 14B as its
  \emph{minimum}. The 0.5B-versus-14B contrast is a contrast with the floor, not with the largest
  model.

\item[{Compile-gain ``$1.34\times$ median''.}]
  Correct and misleading alone. The distribution is bimodal: L1 $15.51\times$, L2 $1.26\times$,
  L3 $4.53\times$. The median is simply L2's value because L2 is the majority tier. Reported per
  tier throughout, because the argument that a plain-torch rewrite scores $0.50$ ``without
  compile being strong'' holds on L2 and fails on L1.
\end{description}

\subsection{Excluded: data that cannot answer the question it was collected for}

\begin{description}[leftmargin=1.4em,style=nextline]
\item[{Severed trajectories (Amendment 5).}]
  \texttt{lab\_compounding} and \texttt{lab\_weight\_rsi} restored round history on resume but not
  the trained adapter. Every resumed chunk continued the round numbering while restarting each
  arm from base weights. The signature is unambiguous: \texttt{trainC} collapses at exactly the
  round named by \texttt{resumed\_at}, in all 18 resumed cells with no exceptions. All affected
  rounds are excluded, and the generator refuses to report a cell whose artifact cannot evidence
  that its weights survived a resume.

\item[{Starved trajectories (Amendment 6).}]
  A trajectory whose self arm received fewer than two verified examples per round on average is
  reported as starved rather than as a measurement of the contrast. At the registered
  configuration every cell qualifies, so \textbf{the registered primary currently reports no
  number at all}, which is the correct outcome rather than a missing one.

\item[{The 32B probe's first reading.}]
  Its first 12 tasks are L1-heavy and gave $6.94\%$. Restricting every scale to those same 12
  tasks is the valid comparison; the full 29-task sweep gives $2.99\%$. The partial figure is
  recorded and not reported.
\end{description}

\subsection{Provisional: retained but not load-bearing}

\begin{description}[leftmargin=1.4em,style=nextline]
\item[{Probe-as-intervention ($+0.124\,[+0.029,+0.218]$).}]
  $n=12$ checkpoints, 5--12 tasks per run, against the weakest available comparator (uniform
  random). Reported in the appendix, not near the abstract.

\item[{T4 closed$\to$open.}]
  $n = 1$ per cell against a registered minimum of three seeds, with no interval. Exploratory.

\item[{Any A100 speed-scored board.}]
  Measured under a harness that silently rejected every non-PyTorch submission. Retained with the
  defect stated, never pooled with H100 sets.
\end{description}

\subsection{Known confounds carried into the results}

\begin{itemize}[leftmargin=1.4em]
\item \textbf{Cross-rung token budgets differ.} T1 and T2 generate at 2048 tokens; T3 and T5 at
  1200. Both resume from per-round checkpoints, so changing them mid-flight would mix budgets
  inside a single cell, which is worse than a consistent difference. Any cross-rung comparison
  carries a $1.7\times$ gap and says so.
\item \textbf{The held set is 5 tasks, not the 20 requested.} A slice does not raise when it runs
  past the end, so a 25-task pool with 20 taken for training returns whatever remains. At $k=10$
  this caps a round at 50 verified candidates, which bounds every per-round score regardless of
  its functional form.
\item \textbf{$n = 2$ trajectories per arm} on the pass-rate board. No interval is quoted on a
  two-point mean; the per-cell values are given instead.
\end{itemize}


\clearpage
\section{Appendix: probe-as-intervention, analysed as a secondary result}
\label{app:probe}
A linear probe reads ``is this kernel correct'' out of mid-to-late layers at high pooled AUC. This appendix asks
whether that decodable signal converts into a \emph{selection} advantage, and reports it at the right unit with
the right denominator.

\paragraph{Three corrections to the earlier analysis.}
(1) \textbf{Unit.} The 19 usable runs are repeated probe-training/test draws over \textbf{12 distinct base
checkpoints}; draws sharing a checkpoint are not independent, so we cluster on checkpoint.
(2) \textbf{Denominator.} Each run's lift is computed over $n_{\text{test}}$ tasks, and
$n_{\text{test}}\in[5,12]$ (median 12). A ``$+0.45$ lift'' is roughly five problems, so
$n_{\text{test}}$ is printed in every row.
(3) \textbf{Estimand.} A high \emph{pooled} correctness AUC can arise because the probe separates easy tasks
from hard ones while failing to rank candidates \emph{within} a task --- and only within-task ranking can
produce a selection gain. Per-candidate probe scores were not retained, so the within-task AUC
\textbf{cannot be recomputed} from the stored artifacts. We flag this as the binding limitation rather than
letting the pooled number stand in for it.

\paragraph{Result.} Clustered on base checkpoint, the mean lift over uniform-random selection is
$+0.124\,[+0.029,+0.218]$ over $n{=}12$ checkpoints --- the 95\% CI \textbf{excludes zero}: against a uniform-random selector at matched generation budget, probe-guided top-1 selection does lift the correct rate.  The result is also \emph{not} TOST-equivalent to zero at the pre-declared $\delta{=}0.05$, so neither a benefit nor its absence would have been concluded by default --- the interval is simply wide. The naive row-level pooling would report
$+0.118\,[+0.028,+0.209]$ over 19 rows. The probe recovers 36\% of the oracle$-$random headroom on average
(95\% CI [11\%,60\%]).

\paragraph{What this does and does not establish.} \textbf{Uniform random is the weakest possible comparator.}
Beating it shows the probe carries \emph{some} usable signal; it does not show the probe is the best use of the
same compute, because the informative comparators were not run: mean/length-normalised log-likelihood reranking,
a compile-or-run validity filter, and execution agreement. Separately, and importantly, this
matched-\emph{generation} selection gain is a different quantity from matched-\emph{verification} harvesting
yield, where we measure no advantage. The coherent joint statement is: \emph{extra draws plus a cheap selector
find more correct kernels, but the tested probe does not make verification-constrained search economically
better.}

\paragraph{The diagnostic.} Several runs achieve pooled $\mathrm{AUC}{=}1.000$ with lifts from $+0.000$ to
$+0.250$, on training sets as small as 2--5 correct examples. Perfect pooled separation that does not predict
selection gain is direct evidence the pooled metric partly measures task difficulty rather than within-task
candidate quality --- a trap worth flagging for anyone reading correctness probes off a small bank.

\paragraph{Defensible claim.} \emph{Decodable correctness information yields a measurable but modest gain over
random selection, and does not automatically translate into useful search efficiency under a verification
budget.} Settling it needs a rerun on a family-split bank of $\geq$300 independent tasks that logs
per-candidate scores, with leave-task-out probe fitting, matched-candidate comparators, task-clustered
uncertainty, and a full generation/scoring/verification cost ledger.

\begin{table}[h]\centering\footnotesize
\caption{Probe-as-intervention, one row per run (not per model). $n_t$ = test tasks the row is computed on;
head\% = fraction of the oracle$-$random headroom recovered. Rows sharing a base checkpoint are repeated draws
and are clustered in the analysis.}
\begin{tabular}{lrrrrrrr}
\toprule
run & $n_t$ & $K$ & AUC & random & probe & lift & head\% \\
\midrule
qwen05k32 & 12 & 32 & 0.831 & 0.218 & 0.667 & +0.449 & 64\% \\
qwen05k32b & 12 & 32 & 0.878 & 0.247 & 0.667 & +0.420 & 63\% \\
qwen05\_nt24 & 5 & 32 & 0.789 & 0.086 & 0.000 & -0.086 & -75\% \\
Qwen1.5\_K32 & 12 & 32 & 0.973 & 0.040 & 0.333 & +0.293 & 64\% \\
qwen15k32 & 11 & 32 & 0.961 & 0.337 & 0.545 & +0.208 & 43\% \\
qwen15\_nt24 & 5 & 32 & 0.806 & 0.050 & 0.200 & +0.150 & 100\% \\
qwen14k16 & 12 & 16 & 0.872 & 0.698 & 0.750 & +0.052 & 17\% \\
Qwen3 & 12 & 16 & 0.958 & 0.266 & 0.667 & +0.401 & 62\% \\
qwen3\_nt24 & 5 & 32 & 0.713 & 0.044 & 0.200 & +0.156 & 100\% \\
qwen3k32b & 12 & 32 & 0.817 & 0.401 & 0.417 & +0.016 & 3\% \\
qwen3\_wt & 12 & 32 & 0.804 & 0.299 & 0.250 & -0.049 & -11\% \\
qwen7\_nt16 & 12 & 32 & 0.824 & 0.531 & 0.667 & +0.136 & 62\% \\
qwen7k32 & 12 & 32 & 0.877 & 0.652 & 0.750 & +0.098 & 28\% \\
qwen7\_nt20 & 9 & 32 & 0.819 & 0.379 & 0.111 & -0.268 & -151\% \\
SmolLM1.7\_K32 & 12 & 32 & 1.000 & 0.083 & 0.333 & +0.250 & 100\% \\
Yi1.5\_K32 & 10 & 32 & 1.000 & 0.100 & 0.100 & +0.000 & -- \\
dsc13k32 & 9 & 32 & 1.000 & 0.178 & 0.222 & +0.044 & 28\% \\
DeepSeek6.7B & 11 & 16 & 0.986 & 0.121 & 0.182 & +0.061 & 40\% \\
dsc67k16 & 12 & 16 & 0.826 & 0.502 & 0.417 & -0.085 & -17\% \\
\bottomrule
\end{tabular}\end{table}


\end{document}
